
\chapter{کارهای پیشین}
\label{chap:literature}

بخش‌بندی خودکار اندام‌های شکمی در پویش بدون‌تزریق دقتی پایین‌تر از پویش دارای ماده‌ی حاجب دارد، در حالی که برای بسیاری از بیماران تنها همین پویش در دسترس است. این فصل ادبیات پژوهشی مرتبط با این مسئله را از دو زاویه مرور می‌کند: نخست و در جایگاه اصلی، هم‌نهادی پویش با ماده‌ی حاجب از روی پویش بدون‌تزریق، که مسئله‌ی سرتاسری این پایان‌نامه است؛ و سپس هم‌ترازسازی تصویر بین‌پویشی، که ساخت جفت‌های آموزشی هم‌نهادی به آن وابسته است.


\section{پیش‌زمینه‌ی بالینی: سی‌تی چندپویشی و پویش‌های تزریق}
\label{sec:clinical-background}

«نوبت تصویربرداری»\footnote{\lr{Study}} یک مراجعه‌ی بیمار برای تصویربرداری است. در سی‌تی شکمی چندپویشی\footnote{\lr{Multi-phase abdominal CT}}، در یک نوبت به‌جای یک پویش واحد، چند پویش متوالی از یک ناحیه‌ی تشریحی، در بازه‌های زمانی مشخص نسبت به تزریق وریدی ماده‌ی حاجب یددار\footnote{\lr{Iodinated contrast medium}}، گرفته می‌شود. پویش نخست، پیش از هرگونه تزریق، پویش بدون‌تزریق نامیده می‌شود و توزیع طبیعی تضعیف اشعه‌ی ایکس\footnote{\lr{X-ray attenuation}} بافت‌ها، بدون اثر ماده‌ی حاجب، را ثبت می‌کند. پس از تزریق، ماده‌ی حاجب به‌ترتیب از قلب به آئورت، سپس به شریان‌های احشایی، و در نهایت از طریق مویرگ‌ها به وریدها و کبد می‌رسد؛ همین توالی زمانی، پایه‌ی نام‌گذاری پویش‌های دارای ماده‌ی حاجب است.

\textbf{پویش شریانی}\footnote{\lr{Arterial phase}}، معمولاً \عدد{۲۵} تا \عدد{۴۰} ثانیه پس از تزریق، زمانی ثبت می‌شود که ماده‌ی حاجب عمدتاً در شریان‌ها و بافت‌های پرخون\footnote{\lr{Hypervascular}} متمرکز است؛ این پویش برای آشکارسازی ضایعاتی مانند کارسینوم سلول کبدی که خون‌رسانی شریانی غالب دارند، حیاتی است. \textbf{پویش وریدی باب}\footnote{\lr{Portal venous phase}}، حدود \عدد{۶۰} تا \عدد{۹۰} ثانیه پس از تزریق، زمانی است که ماده‌ی حاجب پارانشیم کبد را از طریق ورید باب پر کرده و بیشترین تباین میان بافت‌های نرم را فراهم می‌کند؛ این پویش معمول‌ترین پویش تشخیصی برای بخش‌بندی و ارزیابی عمومی اندام‌های شکمی است. \textbf{پویش تأخیری}\footnote{\lr{Delayed phase}}، معمولاً \عدد{۳} تا \عدد{۵} دقیقه پس از تزریق، شست‌وشوی ماده‌ی حاجب از ضایعات پرخون را نشان می‌دهد؛ الگوی تقویت در پویش شریانی و شست‌وشوی آن در پویش تأخیری یا وریدی، معیار تشخیصی اصلی کارسینوم سلول کبدی در دستورعمل‌های بالینی مانند \lr{EASL}~\cite{easl2018} است.


هر یک از این پویش‌ها اطلاعات تشخیصی متفاوت و مکملی را در اختیار رادیولوژیست قرار می‌دهند؛ به همین دلیل، بسیاری از دستورالعمل‌های بالینی به‌طور هم‌زمان چند پویش را ثبت می‌کنند و تفسیر دقیق تنها با ترکیب اطلاعات چند پویش ممکن است. با این حال، همین چندپویشی‌بودن دو چالش محاسباتی به همراه دارد که موضوع اصلی این فصل‌اند: تفاوت شدت روشنایی میان پویش‌ها می‌تواند دقت مدل‌های بخش‌بندی را کاهش دهد (بخش~\ref{sec:seg-gap})، و جابه‌جایی بیمار میان پویش‌های متوالی، هم‌ترازسازی دقیق میان پویش‌ها را ضروری می‌سازد (بخش~\ref{sec:registration}). راه‌حل مکمل، تولید مستقیم پویش با ماده‌ی حاجب از روی پویش بدون‌تزریق است (بخش~\ref{sec:synthesis}).

\section{تعریف مسئله‌ها و معیارهای ارزیابی}
\label{sec:lit-problems}

\subsection*{سه مسئله}

\textbf{شکاف بخش‌بندی.} 
بخش‌بندی چندعضوی تصاویر سی‌تی شکمی، یعنی تعیین دقیق مرزهای تشریحی اندام‌هایی مانند کبد، کلیه، طحال و پانکراس در حجم‌های سه‌بعدی، زیرساخت بسیاری از کاربردهای بالینی مبتنی‌بر هوش مصنوعی از جمله برنامه‌ریزی درمان، اندازه‌گیری حجمی و ردیابی روند بیماری است. این مسئله معمولاً به‌صورت یک وظیفه‌ی دسته‌بندی تصویردانه‌به‌تصویردانه (یا وُکسل‌به‌وُکسل) فرمول‌بندی می‌شود که در آن هر وُکسل حجم سی‌تی باید به یکی از برچسب‌های تشریحی از پیش‌تعریف‌شده نسبت داده‌شود. مشاهده‌ی مبنای این فصل آن است که دقت این بخش‌بندی لزوماً در طول پویش‌های مختلف تصویربرداری یکنواخت باقی نمی‌ماند: پویش بدون‌تزریق، به دلیل کاهش نسبت تباین‌به‌نوفه میان بافت‌های نرم مجاور، اغلب دقت پایین‌تری نسبت به پویش با ماده‌ی حاجب نشان می‌دهد؛ بدون تزریق ماده‌ی حاجب، تفاوت میان بافت نرم اندام‌های مجاور مانند کبد و طحال، یا میان پارانشیم اندام و عروق داخل آن، به‌مراتب کمتر از پویش با ماده‌ی حاجب است و مرزهای تشریحی‌ای که مدل‌های بخش‌بندی برای آن‌ها آموزش دیده‌اند، در تصاویر بدون‌تزریق کم‌رنگ‌تر یا مبهم‌تر ظاهر می‌شوند.

\textbf{هم‌ترازسازی بین‌پویشی.} تراز فضایی میان پویش‌های یک نوبت چندپویشی\footnote{\lr{Multi-phase CT}} پیش‌نیاز بسیاری از کاربردهای بالینی است؛ برای نمونه دستورعمل‌های تشخیص کارسینوم سلول کبدی~\cite{easl2018,forner2018} بر ترکیب اطلاعات پویش شریانی و وریدی تکیه می‌کنند و این ترکیب تنها با انطباق دقیق دو پویش معتبر است. هم‌ترازسازی تصویر به‌طور صوری، مسئله‌ی یافتن یک تبدیل هندسی $T$ است که یک تصویر متحرک\footnote{\lr{Moving image}} را بر یک تصویر ثابت\footnote{\lr{Fixed image}} منطبق کند. چالش مضاعف در هم‌ترازسازی بین‌پویشی آن است که تزریق ماده‌ی حاجب رابطه‌ی آماری میان شدت روشنایی دو تصویر را به‌شکل ناهمگن در فضا تغییر می‌دهد: عروق و پارانشیم پرخون روشن‌تر می‌شوند، در حالی که استخوان، چربی و بسیاری از بافت‌های نرم تقریباً بدون تغییر می‌مانند.

\textbf{هم‌نهادی پویش با ماده‌ی حاجب.} مسیر مکمل، تولید مستقیم اطلاعات پویش با ماده‌ی حاجب از روی تصویر بدون‌تزریق است؛ وظیفه‌ای که با عناوینی مانند \textbf{هم‌نهادی تباین مجازی}\footnote{\lr{Virtual contrast enhancement}} یا ترجمه‌ی تصویر\footnote{\lr{Image-to-image translation}} از \lr{NC} به \lr{CE} شناخته می‌شود و به‌طور صوری یک مسئله‌ی نگاشت شرطی $G: x_{\mathrm{NC}} \mapsto \hat{x}_{\mathrm{CE}}$ است. این حوزه در سال‌های اخیر به‌سرعت رشد کرده و اکنون شامل ده‌ها روش با معماری‌ها، مجموعه‌داده‌ها و راهبردهای نظارتی متفاوت است.

\subsection*{معیارهای بخش‌بندی}


\textbf{ضریب هم‌پوشانی سورنسن--دایس}\footnote{\lr{Sørensen--Dice coefficient}}، یا به‌اختصار \lr{Dice}، معیار غالب در ادبیات بخش‌بندی تصویر پزشکی است. برای دو نقاب دودویی $A$ (پیش‌بینی مدل) و $B$ (برچسب مرجع)، این ضریب به‌صورت $\mathrm{Dice} (A,B) = 2|A\cap B| / (|A|+|B|)$ تعریف می‌شود و مقداری میان صفر (بدون هم‌پوشانی) و یک (هم‌پوشانی کامل) می‌گیرد. مزیت اصلی \lr{Dice} نسبت به دقت تصویردانه‌ای خام، مقاومت نسبی آن در برابر عدم‌توازن حجمی میان پیش‌زمینه (اندام هدف) و پس‌زمینه است، ویژگی‌ای که در بخش‌بندی اندام‌های کوچک مانند پانکراس یا غده‌ی فوق‌کلیوی اهمیت زیادی دارد.

\textbf{صدک نودوپنجم فاصله‌ی هاسدورف}\footnote{\lr{95th percentile Hausdorff Distance (HD95)}} معیاری مکمل است که به خطای موضعی مرزی، برخلاف \lr{Dice}، حساس است. \lr{Dice} تنها میزان هم‌پوشانی حجمی را می‌سنجد و ممکن است خطای برجسته‌ی مرزی در یک ناحیه‌ی کوچک را پنهان کند؛ \lr{HD95} این کاستی را جبران می‌کند: بزرگ‌ترین فاصله‌ی باقی‌مانده میان مرز دو نقاب را، پس از حذف \عدد{۵} درصد دورافتاده‌ترین نقاط (برای کاهش حساسیت به نوفه\footnote{\lr{Noise}} برچسب‌گذاری)، اندازه می‌گیرد. در برخی مطالعات، از انحراف حجمی\footnote{\lr{Volumetric deviation}} نسبت به یک پویش مرجع نیز به‌عنوان معیار مکمل برای سنجش تکرارپذیری\footnote{\lr{Reproducibility}} استفاده می‌شود.

\subsection*{معیارهای هم‌ترازسازی}


\textbf{اطلاعات متقابل}\footnote{\lr{Mutual Information (MI)}} و نسخه‌ی بهنجارشده‌ی آن\footnote{\lr{Normalised Mutual Information (NMI)}} معیارهای آماری‌اند که وابستگی مشترک توزیع شدت روشنایی دو تصویر را می‌سنجند، بدون فرض رابطه‌ی خطی میان آن‌ها. این معیارها هم به‌عنوان تابع هزینه‌ی بهینه‌سازی در حین هم‌ترازسازی و هم به‌عنوان شاخص کیفیت پس از هم‌ترازسازی به کار می‌روند، و پایه‌ی معیار تراز در بسیاری از روش‌های کلاسیک (\cite{wells1996,studholme1999,mattes2001}) و ابزارهایی مانند \lr{elastix} و \lr{SimpleITK}اند.

\textbf{همبستگی متقابل نرمال‌شده}\footnote{\lr{Normalised Cross-Correlation (NCC)}} شباهت خطی شدت روشنایی یا شیو\footnote{\lr{Gradient}} دو تصویر را می‌سنجد. این معیار پایه‌ی \lr{SyN}~\cite{avants2008} و توصیف‌گرهای ساختاری مانند \lr{MIND}~\cite{heinrich2012} و \lr{DEEDS}~\cite{heinrich2013} است.

\textbf{خطای هم‌ترازی هدف (TRE) }\footnote{\lr{Target Registration Error (TRE)}} میانگین فاصله‌ی باقی‌مانده میان نشانه‌های آناتومیک متناظر\footnote{\lr{Anatomical landmarks}} پس از هم‌ترازسازی را اندازه می‌گیرد؛ این معیار مستقیم‌ترین سنجه‌ی دقت هندسی است و در محک‌هایی مانند \lr{Learn2Reg}~\cite{murphy2011} برای گزارش دقت روش‌هایی مانند \lr{DEEDS} و \lr{TransMorph} به کار می‌رود.

هر دو خانواده‌ی معیار محدودیت دارند. \lr{TRE} به نشانه‌گذاری دستی پرهزینه نیاز دارد و تنها برای زیرمجموعه‌ی کوچکی از داده‌ها در دسترس است؛ به همین دلیل بسیاری از مطالعات، از جمله \lr{VoxelMorph}~\cite{balakrishnan2019}، \lr{TransMorph}~\cite{chen2022} و \lr{DEEDS}~\cite{heinrich2013}، از هم‌پوشانی \lr{Dice} نقاب‌های تشریحی به‌عنوان جایگزین استفاده می‌کنند. این جایگزینی نیز دو محدودیت دارد: \lr{Dice} تنها هم‌پوشانی ناحیه‌ای را می‌سنجد و نه تناظر نقطه‌به‌نقطه، و کیفیتش به دقت نقاب‌های ورودی (معمولاً حاصل یک مدل بخش‌بندی خودکار) وابسته است. معیارهای شباهت شدت (\lr{MI}، \lr{NCC}) نیز تنها پروکسی غیرمستقیم تراز هندسی‌اند و در نواحی با تغییر ناهمگن تباین ممکن است حتی بدون تراز واقعی مقدار گمراه‌کننده‌ای گزارش کنند.

\subsection*{معیارهای هم‌نهادی}


ارزیابی کیفیت هم‌نهادی معمولاً هم به وفاداری تصویردانه‌ای و هم به واقع‌نمایی ادراکی تصویر تولیدشده توجه می‌کند.

\textbf{نسبت علامت\footnote{\lr{Signal}} به نوفه اوج}\footnote{\lr{Peak Signal-to-Noise Ratio (PSNR)}} بر پایه‌ی میانگین مربعات خطا در فضای لگاریتمی تعریف می‌شود و تصویر هم‌نهاده\footnote{\lr{Synthetic}} را مستقیماً با تصویر واقعی مقایسه می‌کند؛ این معیار بیشتر به خطای تصویردانه‌ای خام حساس است تا ساختار ادراکی تصویر.

\textbf{شاخص شباهت ساختاری}\footnote{\lr{Structural Similarity Index (SSIM)}} نیز تصویر هم‌نهاده را مستقیماً با تصویر واقعی مقایسه می‌کند، اما با سنجش هم‌زمان روشنایی، تباین و ساختار محلی، به تغییرات ادراکی نزدیک‌تر از \lr{PSNR} است.

\textbf{میانگین قدرمطلق خطا}\footnote{\lr{Mean Absolute Error (MAE)}}، \textbf{میانگین مربعات خطا}\footnote{\lr{Mean Squared Error (MSE)}} و \textbf{ضریب همبستگی پیرسون}\footnote{\lr{Pearson Correlation Coefficient (PCC)}} معیارهای متعارف تفاوت تصویردانه‌به‌تصویردانه‌اند؛ دو مورد نخست در بسیاری از مطالعات بخشی از تابع هزینه‌ی آموزش نیز هستند.

\textbf{فاصله‌ی آغازشی فرشه}\footnote{\lr{Fréchet Inception Distance (FID)}} برخلاف معیارهای فوق، تصویر به تصویر مقایسه نمی‌کند، بلکه فاصله‌ی میان توزیع ویژگی‌های استخراج‌شده (با یک شبکه‌ی از پیش‌آموزش‌دیده) از کل مجموعه‌ی تصاویر واقعی و تولیدشده را اندازه می‌گیرد؛ پس معیاری در سطح توزیع، نه تک‌تصویر، است.

\textbf{شباهت پرسپچوال یادگرفته‌شده}\footnote{\lr{Learned Perceptual Image Patch Similarity (LPIPS)}} تفاوت دو تصویر را در فضای ویژگی یک شبکه‌ی از پیش‌آموزش‌دیده، به‌جای فضای تصویردانه\footnote{\lr{Pixel}} خام، می‌سنجد و بنابراین با قضاوت ادراکی انسان هم‌بستگی بهتری نسبت به \lr{PSNR}/\lr{SSIM} دارد.

سرانجام، در مطالعاتی که کاربردپذیری پایین‌دستی هم‌نهادی را می‌سنجند، ضریب \textbf{\lr{Dice}} بخش‌بندی روی تصویر هم‌نهاده (تعریف‌شده در همین بخش) به‌عنوان معیار غیرمستقیم کیفیت به کار می‌رود.

\section{شکاف بخش‌بندی در پویش بدون‌تزریق}
\label{sec:seg-gap}


چارچوب \textbf{\lr{nnU-Net}}~\cite{isensee2021nnunet} یک شبکه‌ی رمزگذار--رمزگشای پیچشی بر پایه‌ی معماری \lr{U-Net} است که با \textbf{خودپیکربندی}\footnote{\lr{Self-configuring}} خودکار معماری، پیش‌پردازش و آموزش بر اساس ویژگی‌های هر مجموعه‌داده، نیاز به تنظیم دستی فراپراسنجه\footnote{\lr{Hyperparameter}}ها را حذف می‌کند و به یکی از پرکاربردترین مبناهای این حوزه بدل شده است. این چارچوب روی \عدد{۲۳} مجموعه‌داده‌ی عمومی از چالش‌های بین‌المللی ارزیابی شده است؛ در وظیفه‌های \emph{شکمی} محک \lr{MSD}~\cite{msd2022}، \lr{Dice} آن برای طحال \عدد{۰٫۹۶}، کبد \عدد{۰٫۹۳} و پانکراس \عدد{۰٫۷۹} گزارش شده است، و برای ساختارهای دشوارتر همان محک به‌مراتب پایین‌تر می‌آید: عروق کبدی \عدد{۰٫۶۳}، تومور کبد \عدد{۰٫۷۴}، توده‌ی پانکراس \عدد{۰٫۵۲} و سرطان روده‌ی بزرگ \عدد{۰٫۵۶}. همه‌ی این ارقام روی داده‌ی پویش با ماده‌ی حاجب به دست آمده‌اند؛ محک \lr{MSD} هیچ وظیفه‌ی بدون‌تزریقی ندارد.

\textbf{\lr{TotalSegmentator}}~\cite{wasserthal2023} یک مدل \lr{nnU-Net} سه‌بعدی است که بر پایه‌ی چارچوب فوق و با آموزش روی \عدد{۱٬۲۰۴} نوبت سی‌تی بالینی متنوع ساخته شده و \عدد{۱۰۴} ساختار تشریحی از جمله اندام‌ها، استخوان‌ها، عضلات و عروق را بخش‌بندی می‌کند؛ این مدل \lr{Dice} برابر \عدد{۰٫۹۴۳} روی مجموعه‌ی آزمون خود گزارش کرده و در عمل به ابزاری مرجع برای تولید نقاب\footnote{\lr{Mask}} اندام در پژوهش‌های بعدی تبدیل شده است. این رقم روی مجموعه‌ی آزمونی گزارش شده که از نظر پویش تزریق \emph{ترکیبی} است و مقاله هیچ تفکیک پویشی از عملکرد ارائه نمی‌کند.

یو و همکاران~\cite{yu2022multiorgan} با یک \lr{nnU-Net} سه‌بعدی، هشت ساختار شکمی را روی تصاویر بدون‌تزریق\footnote{\lr{Non-contrast CT (NCCT)}} و با ماده‌ی حاجب\footnote{\lr{Contrast-enhanced CT (CECT)}} بخش‌بندی کردند. نکته‌ی روش‌شناختی مهم در این مطالعه آن است که ضریب \lr{Dice} \emph{تنها} برای پویش با ماده‌ی حاجب گزارش شده است (از \عدد{۰٫۸۰} برای دوازدهه تا \عدد{۰٫۹۷} برای طحال، روی \عدد{۳۰} بیمار آزمون)، و مقایسه‌ی میان دو پویش صرفاً به‌صورت \emph{کیفی} و با امتیازدهی لیکرت رادیولوژیست انجام شده است: \عدد{۸۹٫۳} درصد خطوط دوازدهه در پویش با ماده‌ی حاجب در برابر \عدد{۸۵٫۳} درصد در پویش بدون‌تزریق، امتیاز سه یا بالاتر (نیازمند ویرایش جزئی) گرفتند. بنابراین این مطالعه، برخلاف آن‌چه از عنوانش برمی‌آید، شواهد کمّی مستقیمی درباره‌ی اندازه‌ی شکاف بین‌پویشی به دست نمی‌دهد.

جئون و همکاران~\cite{jeon2024fully} تنها مطالعه‌ی این مجموعه‌اند که شکاف بین‌پویشی را به‌صورت کمّی و روی داده‌ی کنترل‌شده اندازه گرفته‌اند. آن‌ها با سی‌تی دوانرژی\footnote{\lr{Dual-energy CT}} جفت‌های ذاتاً هم‌ترازِ پویش وریدی و تصویر بدون‌تزریق مجازی\footnote{\lr{Virtual non-contrast (VNC)}} ساختند، یک \lr{nnU-Net} سه‌بعدی را روی \عدد{۷۵} نوبت آموزش دادند و آن را روی \عدد{۳۰} نوبت سی‌تی تک‌انرژی خارجی آزمودند. الگوی نتیجه روشن است: برای اندام‌های بزرگ و پرتباین (کبد، طحال، دو کلیه)، دقت از پویش وریدی به پویش بدون‌تزریق تنها اندکی افت می‌کند، اما برای پانکراس، اندامی کوچک با مرز مبهم در برابر بافت چرب اطراف، این افت به‌روشنی بزرگ‌تر است. ارقام کامل \lr{Dice} در هر دو پویش در جدول~\ref{tab:seggap-summary} آمده است.

پارک و همکاران~\cite{park2025nect}، از همان گروه پژوهشی جئون، مسیر دیگری را آزمودند: به‌جای ساختن جفت هم‌تراز، مدل را \emph{مستقیماً و تنها} روی داده‌ی بدون‌تزریق آموزش دادند. داده‌ی توسعه \عدد{۱۰۵} نوبت سی‌تی شکمی بدون‌تزریق از بخش اورژانس یک بیمارستان دانشگاهی و \عدد{۶۰} سی‌تی کم‌دوز قفسه‌سینه از برنامه‌ی غربالگری همان مرکز بود. نتیجه چشمگیر است: \lr{Dice} روی مجموعه‌ی آزمون داخلی برای کبد \عدد{۰٫۹۸۲}، طحال \عدد{۰٫۹۷۳}، کلیه‌ی راست \عدد{۰٫۹۶۶} و کلیه‌ی چپ \عدد{۰٫۹۶۷}، و روی مجموعه‌ی آزمون دوم به‌ترتیب \عدد{۰٫۹۶۸}، \عدد{۰٫۹۶۰}، \عدد{۰٫۹۵۲} و \عدد{۰٫۹۵۸} بود؛ یعنی هم‌سطح با ارقامی که مدل‌های دیگر روی پویش \emph{با ماده‌ی حاجب} گزارش می‌کنند. اختلاف حجم برآوردشده با حجم مرجع نیز میان $-0.7$ تا $+2.2$ درصد ماند.

دامنه‌ی تعمیم این ارقام اما سه قید دارد که هر سه از خودِ مقاله برمی‌آیند. نخست، آن‌چه مقاله «مجموعه‌ی آزمون خارجی» می‌نامد، \عدد{۵۵} نوبت شکمی و \عدد{۱۰} نوبت قفسه‌سینه از \emph{همان بیمارستان} است که تنها با پویش‌گر\footnote{\lr{Scanner}}های متفاوت گرفته شده‌اند؛ پس این ارزیابی، جابه‌جایی حوزه‌ی ناشی از تغییر دستگاه را می‌سنجد و نه تغییر مرکز، جمعیت یا دستورالعمل. دوم، فهرست اندام‌ها به چهار اندام توپر و بزرگ محدود است و پانکراس، یعنی همان اندامی که در مطالعه‌ی جئون بیشترین افت بین‌پویشی (\عدد{۰٫۰۳۶}) را نشان داد، در آن نیست. سوم، معیارهای ورود، بیماران دارای ضایعه‌ی تغییردهنده‌ی کانتور را کنار گذاشته و قرارداد برچسب‌گذاری نیز مرزهای مبهم بدون ماده‌ی حاجب (عروق درون کبد، سینوس کلیه) را از پیش از نقاب بیرون گذاشته است؛ پس عدد گزارش‌شده، سازگاری با یک قرارداد خودسازگار را روی جمعیتی نزدیک‌به‌طبیعی می‌سنجد. خود نویسندگان نیز در محدودیت‌ها به کوچک‌بودن مجموعه‌ی آزمون و طبیعی‌بودن اندام‌ها تصریح کرده‌اند.

با این قیدها، این مطالعه یک نکته‌ی مهم را تثبیت می‌کند: پایین‌بودن دقت در پویش بدون‌تزریق یک محدودیت فیزیکی و اجتناب‌ناپذیر نیست. دست‌کم برای اندام‌های بزرگ توپر و در جمعیتی همگن، مدلی که مستقیماً روی داده‌ی بدون‌تزریق آموزش ببیند به همان سطحی می‌رسد که مدل‌های آموزش‌دیده روی پویش دارای ماده‌ی حاجب، روی داده‌ی خودشان، گزارش می‌کنند. اما راه‌حل پیشنهادی آن مقیاس‌پذیر نیست: به داده‌ی بدون‌تزریقِ برچسب‌خورده و بومیِ هر مرکز و هر مجموعه‌ی اندام نیاز دارد، و فصل~\ref{chap:survey} مستند می‌کند که چنین داده‌ای به‌اندازه‌ی کافی وجود ندارد و برچسب‌گذاری دستی آن به‌مراتب دشوارتر از پویش با ماده‌ی حاجب است. همین محدودیت است که در ادبیات، دو مسیر جایگزین را پیش کشیده است: هم‌ترازسازی پویش‌های موجود (بخش~\ref{sec:registration}) و هم‌نهادی پویش غایب (بخش~\ref{sec:synthesis}).

آبل و همکاران~\cite{abel2025intraindividual} به‌جای دقت مطلق، \emph{ثبات} خروجی را سنجیدند: آن‌ها \عدد{۳٬۷۵۶} سری سی‌تی را که هر بیمار در سه پویش شریانی، وریدی و بدون‌تزریق داشت به \lr{TotalSegmentator} دادند و انحراف حجمی هر ساختار را میان پویش‌ها اندازه گرفتند. \عدد{۲۹} ساختار از \عدد{۳۴} ساختار، انحراف حجمی کمتر از \عدد{۵} درصد نشان دادند و پس تکرارپذیر تلقی شدند؛ تنها اندام‌های معدودی مانند غده‌ی فوق‌کلیوی، کیسه‌ی صفرا، طحال و دوازدهه به \textbf{جابه‌جایی حوزه‌ای}\footnote{\lr{Domain shift}} وابسته به پویش و دستورالعمل حساس‌تر بودند. این مطالعه دو نکته را برجسته می‌کند: نخست آن‌که تغییر پویش، خروجی یک ابزار ثابت را به‌طور نظام‌مند جابه‌جا می‌کند؛ و دوم آن‌که این جابه‌جایی، مانند یافته‌ی جئون، روی زیرمجموعه‌ی کوچکی از ساختارها متمرکز است.

ژو و همکاران~\cite{Zho24} چارچوب \textbf{\lr{CKAP-Net}} را پیشنهاد دادند که یک پیمانه\footnote{\lr{Module}} تزریق دانش تباین و یک شاخه‌ی ادراک تشریحی چندمقیاسی را برای هم‌نهادی اطلاعات تباین از روی تصویر بدون‌تزریق، پیش از بخش‌بندی، ترکیب می‌کند؛ این مدل به ضریب \lr{Dice} حدود \عدد{۰٫۸۱۲} در بخش‌بندی چندعضوی روی مجموعه‌داده‌های چندمرکزی رسید، رقمی که فاصله‌ی چندانی با دقت متداول بخش‌بندی روی پویش با ماده‌ی حاجب ندارد.

جدول~\ref{tab:seggap-summary} این مطالعات را کنار هم می‌گذارد. مطالعه‌ی آبل و همکاران به‌دلیل آن‌که معیارش \lr{Dice} نیست در جدول نیامده است. خالی‌بودن گسترده‌ی ستون \lr{Dice} پویش بدون‌تزریق خود مهم‌ترین یافته‌ی این بخش است: هیچ‌یک از مدل‌های مرجع عملکرد خود را به‌تفکیک پویش گزارش نکرده‌اند، و تنها سه مطالعه (جئون، پارک و \lr{CKAP-Net}) برای پویش بدون‌تزریق عددی دارند که هر سه برای به‌دست‌آوردنش ناچار شده‌اند داده‌ی اختصاصی بسازند. ستون «داده‌ی آموزش» به همین دلیل افزوده شده است: بدون دانستن منشأ داده‌ی هر مدل، مقایسه‌ی دو ستون \lr{Dice} معنا ندارد.

شکاف بین‌پویشی \emph{وجود دارد} اما یکنواخت نیست: در تنها اندازه‌گیری کمّی و کنترل‌شده‌ی موجود~\cite{jeon2024fully}، افت روی اندام‌های بزرگ و پرتباین میان \عدد{۰٫۰۱۱} و \عدد{۰٫۰۲۷} و روی اندام کوچک و کم‌مرزی مانند پانکراس \عدد{۰٫۰۳۶} است. شکاف \emph{تماماً ذاتی نیست}: نتیجه‌ی پارک و همکاران~\cite{park2025nect} نشان می‌دهد که برای اندام‌های بزرگ توپر و در یک جمعیت همگن و نزدیک‌به‌طبیعی، آموزش مستقیم روی داده‌ی بدون‌تزریق بخش عمده‌ی آن را می‌بندد؛ اما چنان‌که در بالا آمد، این شاهد نه اندام‌های دشوار را پوشش می‌دهد و نه عبور از مرز مرکز را، و بنابراین آن را باید \emph{کران بالای} آن‌چه با داده‌ی متناسب به دست می‌آید خواند، نه توصیفی از حالت کلی. و سرانجام، همین راه‌حل به محدودیت داده برمی‌خورد؛ ادبیات پژوهشی دو راهبرد مکمل برای جبران این افت پیشنهاد داده است که در بخش‌های~\ref{sec:registration} و~\ref{sec:synthesis} همین فصل مرور می‌شوند: هم‌ترازسازی هندسی پویش بدون‌تزریق با پویش دارای ماده‌ی حاجب، و هم‌نهادی مستقیم اطلاعات تباین از روی تصویر بدون‌تزریق.

{\scriptsize
\begin{longtable}{|p{1.9cm}|p{0.45cm}|p{1.85cm}|p{2.5cm}|p{1.15cm}|p{1.5cm}|p{1.7cm}|p{0.45cm}|}
\caption{مدل‌های بخش‌بندی چندعضوی و مطالعات شکاف بین‌پویشی.
\textbf{کد پویش:} \عدد{۰} بدون‌تزریق، \عدد{۱} شریانی، \عدد{۲} وریدی باب، \عدد{۳} تأخیری؛ ستون «پویش‌های داده‌ی آموزش» فهرست پویش‌هایی است که مدل در آموزش دیده است.
همه‌ی مقادیر \lr{Dice} در بازه‌ی صفر تا یک‌اند. \textbf{ستون کد:} \checkmark یعنی کد منبع به‌طور تأییدشده در دسترس عمومی است و \xmark یعنی کد منتشر نشده یا وضعیت آن در مقاله ذکر نشده است. جزئیات و قیدهای هر سطر در متن بالا آمده است.}
\label{tab:seggap-summary} \\
\hline
\textbf{روش (مرجع)} & \textbf{سال} & \textbf{اندام‌های هدف} & \textbf{نوع و منشأ داده‌ی آموزش} & \textbf{پویش‌های آموزش} & \textbf{\lr{Dice} پویش با ماده‌ی حاجب} & \textbf{\lr{Dice} پویش بدون‌تزریق} & \textbf{کد} \\
\hline
\endfirsthead
\multicolumn{8}{r}{\scriptsize \emph{ادامه‌ی جدول~\ref{tab:seggap-summary} از صفحه‌ی پیش}} \\
\hline
\textbf{روش (مرجع)} & \textbf{سال} & \textbf{اندام‌های هدف} & \textbf{نوع و منشأ داده‌ی آموزش} & \textbf{پویش‌های آموزش} & \textbf{\lr{Dice} پویش با ماده‌ی حاجب} & \textbf{\lr{Dice} پویش بدون‌تزریق} & \textbf{کد} \\
\hline
\endhead
\hline
\multicolumn{8}{r}{\scriptsize \emph{ادامه در صفحه‌ی بعد}} \\
\endfoot
\endlastfoot
\lr{nnU-Net}~\cite{isensee2021nnunet} & \عدد{۲۰۲۱} & اندام و تومور شکمی (وظیفه‌های شکمی \lr{MSD}) & \عدد{۲۳} مجموعه‌داده‌ی عمومی چالش‌ها؛ سی‌تی با ماده‌ی حاجب چندمرکزی & \عدد{۲} & اندام: \عدد{۰٫۷۹--۰٫۹۶} \newline تومور: \عدد{۰٫۵۲--۰٫۷۴} & --- & \checkmark \\
\hline
\lr{Total\-Segmen\-tator}~\cite{wasserthal2023} & \عدد{۲۰۲۳} & \عدد{۱۰۴} ساختار تشریحی کل‌بدن & \عدد{۱٬۲۰۴} نوبت بالینی روتین یک مرکز؛ پویش‌گر و دستورالعمل متنوع & \عدد{۰}، \عدد{۱}، \newline \عدد{۲}، \عدد{۳} \newline (تفکیک‌نشده) & \عدد{۰٫۹۴۳} (ترکیبی) & --- & \checkmark \\
\hline
\lr{CLIP-driven Universal Model}~\cite{liu2023clipdriven} & \عدد{۲۰۲۳} & \عدد{۲۵} اندام + شش نوع تومور & \عدد{۳٬۴۱۰} نوبت از \عدد{۱۴} مجموعه‌داده‌ی عمومی؛ با ماده‌ی حاجب & \عدد{۲} & \عدد{۰٫۸۶۱} (\lr{BTCV}) & --- & \checkmark \\
\hline
یو و همکاران~\cite{yu2022multiorgan} & \عدد{۲۰۲۲} & هشت ساختار شکمی، شامل دوازدهه & درون‌سازمانی؛ هر بیمار در دو پویش، \عدد{۳۰} بیمار آزمون & \عدد{۰} و \عدد{۲} & \عدد{۰٫۸۰--۰٫۹۷} & --- & \xmark \\
\hline
جئون و همکاران~\cite{jeon2024fully} & \عدد{۲۰۲۴} & کبد، طحال، دو کلیه، پانکراس & \عدد{۷۵} نوبت سی‌تی دوانرژی؛ جفت ذاتاً هم‌ترازِ وریدی و بدون‌تزریق مجازی. آزمون: \عدد{۳۰} نوبت تک‌انرژی خارجی & \عدد{۰} (مجازی) و \عدد{۲} & تک‌انرژی: \عدد{۰٫۸۴۶--۰٫۹۸۱} & مجازی: \عدد{۰٫۸۵۴--۰٫۹۷۷} \newline واقعی: \عدد{۰٫۸۱۰--۰٫۹۶۵} & \xmark \\
\hline
پارک و همکاران~\cite{park2025nect} & \عدد{۲۰۲۵} & کبد، طحال، دو کلیه (بدون پانکراس) & \عدد{۱۰۵} سی‌تی شکمی بدون‌تزریق اورژانس + \عدد{۶۰} سی‌تی کم‌دوز قفسه‌سینه؛ یک مرکز، آناتومی نزدیک‌به‌طبیعی & \textbf{\عدد{۰} تنها} & --- & شکمی: \عدد{۰٫۹۵۲--۰٫۹۶۸} \newline قفسه‌سینه: \عدد{۰٫۹۶۰--۰٫۹۶۹} & \xmark \\
\hline
\lr{CKAP-Net}~\cite{Zho24} & \عدد{۲۰۲۴} & \عدد{۱۳} اندام شکمی (هم‌نهادی + بخش‌بندی) & \عدد{۴۰۳} و \عدد{۳۸} بیمار از دو مرکز؛ جفت‌های بدون‌تزریق و شریانی دیررس & \عدد{۰} و \عدد{۱} (جفت‌شده) & --- & \عدد{۰٫۸۱۲} & \xmark \\
\hline
\end{longtable}
}


\section{هم‌ترازسازی بین‌پویشی سی‌تی}
\label{sec:registration}

این بخش هم‌ترازسازی را در حدی مرور می‌کند که برای فهم روش پیشنهادی لازم است: خطوط پژوهشی عمومی، سپس تنها مطالعه‌ی هم‌ترازسازی بین‌پویشی شکمی، و در پایان جای خالی.

\subsection*{معیارهای شدت‌محور و مدل‌های تغییر شکل کلاسیک}


پایه‌ی نظری هم‌ترازسازی خودکار چندوجهی را ولز و همکاران~\cite{wells1996} با معرفی \textbf{اطلاعات متقابل} گذاشتند. مسئله این بود که وقتی دو تصویر از دو وجهیت متفاوت می‌آیند، رابطه‌ی میان شدت روشنایی متناظر در آن‌ها نه خطی است و نه از پیش معلوم، و معیارهایی مانند مجموع مربعات تفاوت بی‌معنا می‌شوند. اطلاعات متقابل، که میزان کاهش عدم‌قطعیت درباره‌ی یک تصویر را در صورت دانستن دیگری می‌سنجد، هیچ مدلی برای رابطه‌ی شدت‌ها فرض نمی‌کند؛ تنها فرض می‌کند چنین رابطه‌ای، هرچه باشد، \emph{در سراسر تصویر یکسان} است. همین فرض ظاهراً بی‌ضرر، همان‌جایی است که در سی‌تی چندپویشی می‌شکند.

دو اصلاح بعدی این معیار را به شکل امروزی‌اش رساندند: استادهولم و همکاران~\cite{studholme1999} با \textbf{اطلاعات متقابل بهنجارشده} وابستگی آن به اندازه‌ی ناحیه‌ی هم‌پوشانی را برطرف کردند، و مَتیس و همکاران~\cite{mattes2001} با هموارسازی بافت‌نگاشت\footnote{\lr{Histogram}} مشترک (\textbf{پنجره‌ی پارزن}) صورتی مشتق‌پذیر از آن ساختند که هنوز هسته‌ی \lr{elastix}~\cite{klein2010} و \lr{SimpleITK}~\cite{lowekamp2013} است. موازی با آن، مدل‌سازی خود تبدیل هندسی مسیر جداگانه‌ای پیمود: الگوریتم \textbf{دیمانز}~\cite{thirion1998} حجم‌دانه\footnote{\lr{Voxel}}‌ها را بر اساس شیو محلی شدت روشنایی می‌راند، ساده و سریع اما شکننده در برابر تغییر تباین؛ \textbf{تغییر شکل آزاد بی‌اسپلاین}~\cite{rueckert1999} جابه‌جایی را با شبکه‌ای تُنُک از نقاط کنترل پراسنجه‌ای می‌کند و هموار‌بودن آن را ذاتاً تضمین می‌کند؛ و \lr{SyN}~\cite{avants2008} با تقارن و دیفئومورفیسم، هسته‌ی جعبه‌ابزار \lr{ANTs} شد.

\subsection*{چرا این خانواده در سی‌تی چندپویشی می‌شکند}


مشکل بنیادین، همان فرضی است که در آغاز این بخش برجسته شد: خانواده‌ی اطلاعات متقابل فرض می‌کند رابطه‌ی آماری میان شدت روشنایی دو تصویر در سراسر میدان دید \emph{پایدار} است. تزریق ماده‌ی حاجب دقیقاً همین فرض را نقض می‌کند، و آن را به‌شکلی نقض می‌کند که بدترین حالت ممکن است: تغییر، \emph{ناهمگن در فضا} است. آئورت و ورید باب در پویش شریانی صدها واحد \lr{HU} روشن‌تر می‌شوند، پارانشیم کبد چند ده واحد، و استخوان و چربی تقریباً هیچ. نتیجه آن است که معیار تراز کلی‌نگر، بخش بزرگی از وردایی\footnote{\lr{Variance}} خود را از تغییر تباین می‌گیرد، نه از جابه‌جایی هندسی واقعی. بهینه‌ساز، تصویر را تغییر شکل می‌دهد تا الگوی \emph{شدت} را تطبیق دهد، در حالی که هدف مطلوب تطبیق \emph{آناتومی} است.

\subsection*{توصیف‌گرهای ساختاری و روش‌های یادگیری‌محور}


مسیر مفهومی متفاوتی برای دور زدن ناپایداری شدت روشنایی وجود دارد: به‌جای مقایسه‌ی \emph{مقدار} شدت، \emph{الگوی محلی} آن مقایسه شود. \textbf{\lr{MIND}}~\cite{heinrich2012} این ایده را عملیاتی کرد. برای هر حجم‌دانه، به‌جای مقدار خام، برداری محاسبه می‌شود که فاصله‌ی وصله‌ی پیرامون آن حجم‌دانه را با وصله‌های همسایه‌اش کد می‌کند؛ یعنی \textbf{خودشباهت} محلی. اگر کل تصویر با یک تبدیل یکنوای شدت روشنایی تغییر کند، این الگوی نسبی تقریباً دست‌نخورده می‌ماند. توصیف‌گر حاصل، شباهتی مستقل از وجهیت فراهم می‌کند و مسئله را از «تطبیق شدت‌ها» به «تطبیق ساختارها» منتقل می‌کند.

بر همین پایه، \textbf{\lr{DEEDS}}~\cite{heinrich2013} توصیف‌گر \lr{MIND} را با \textbf{بهینه‌سازی گسسته‌ی میدان تصادفی مارکوف} ترکیب کرد. به‌جای بهینه‌سازی پیوسته و مبتنی‌بر شیو (که در کمینه‌های محلی گیر می‌کند) فضای جابه‌جایی گسسته‌سازی می‌شود و بهترین ترکیب جابه‌جایی‌ها با انتقال پیام روی یک گراف یافت می‌شود. این ترکیب، \lr{DEEDS} را به یکی از دقیق‌ترین روش‌های هم‌ترازسازی شکمی در محک \lr{Learn2Reg} بدل کرده است: \lr{Dice} برابر \عدد{۰٫۷۷۵} در هم‌ترازسازی درون‌وجهی \lr{MR} و \عدد{۰٫۶۳۶} در هم‌ترازسازی بین‌وجهی \lr{MR}--\lr{CT}. افت چشمگیر میان این دو رقم نشان می‌دهد حتی توصیف‌گرهای ساختاری نیز در برابر تغییر وجهیت کاملاً مصون نیستند.

روش‌های یادگیری‌محور به‌جای بهینه‌سازی تکراری برای هر جفت، یک شبکه را یک‌بار آموزش می‌دهند تا میدان جابه‌جایی را در یک گذر پیش‌رو پیش‌بینی کند. \lr{VoxelMorph}~\cite{balakrishnan2019} الگوی پایه‌ی این خانواده را تثبیت کرد (\lr{Dice} برابر \عدد{۰٫۷۵۳} بدون‌نظارت و \عدد{۰٫۷۸} با نظارت، روی هشت محک مغزی)، و \lr{TransMorph}~\cite{chen2022}، \lr{LapIRN}~\cite{mok2020}، شبکه‌های آبشاری~\cite{zhao2019} و \lr{HyperMorph}~\cite{hoopes2021} هر یک یکی از محدودیت‌های آن را هدف گرفتند. \textbf{\lr{UniGradICON}}~\cite{tian2023}، تازه‌ترین گام، یک \textbf{مدل بنیادین} هم‌ترازسازی است که روی دوازده مجموعه‌داده‌ی چندآناتومی آموزش دیده است؛ عملکرد آن به‌شدت به آناتومی وابسته است: \lr{Dice} برابر \عدد{۰٫۷۶۲} روی مغز اما تنها \عدد{۰٫۴۸۳} روی شکم، و رفتارش در برابر جابه‌جایی حوزه‌ای ناشی از \emph{پویش تزریق} هنوز مشخصه‌یابی نشده است.

\subsection*{روش‌های آگاه از آناتومی}

چند خط پژوهشی کوشیده‌اند اطلاعات تشریحی را وارد مسئله کنند. پاپیه‌ژ و همکاران~\cite{papiez2014} با یک منظم‌ساز آگاه از آناتومی حرکت لغزشی در مرز اندام‌ها را ممکن کردند، و دِ‌فوس و همکاران~\cite{devos2019} سهم هر ناحیه در تابع هزینه‌ی یک شبکه‌ی پیچشی را بر اساس نوع بافت وزن‌دهی کردند. هو و همکاران~\cite{hu2018} با \textbf{نظارت ضعیف}\footnote{\lr{Weak supervision}} مبتنی‌بر هم‌پوشانی برچسب، تابع هزینه را به‌کلی از شدت روشنایی جدا کردند، به بهای نیاز به پیکره‌ای از جفت‌های \emph{برچسب‌خورده} و نگاشتی آموخته‌شده که بدون بازآموزی به مرکز یا دستورالعمل جدید منتقل نمی‌شود. اوکتای و همکاران~\cite{oktay2017} نیز \textbf{پیشین‌های شکل تشریحی}\footnote{\lr{Anatomical shape priors}} را به‌عنوان منظم‌ساز به تابع هزینه افزودند.


نقطه‌ی مشترک این خط پژوهشی آن است که همگی از آناتومی برای \emph{وزن‌دهی} به معیار تراز استفاده می‌کنند، نه برای حذف نواحی ناپایدار از \emph{محاسبه‌ی} آن. این تفاوت در نگاه اول جزئی به نظر می‌رسد، اما نیست. اطلاعات متقابل از یک بافت‌نگاشت مشترک ساخته می‌شود که از همه‌ی حجم‌دانه‌ها تغذیه می‌کند. وزن کم دادن به روده تنها سهم آن را در جمع نهایی کم می‌کند؛ حجم‌دانه‌های روده همچنان در ساخت بافت‌نگاشت شرکت می‌کنند و آن را چندمُدی نگه می‌دارند. حذف کامل روده از دامنه‌ی محاسبه، برخلاف وزن‌دهی، خودِ بافت‌نگاشت را عوض می‌کند.

\subsection*{هم‌ترازسازی چندپویشی شکمی}


برخلاف آن‌چه از حجم ادبیات بالا برمی‌آید، مسئله‌ی مشخصِ هم‌ترازسازی \emph{بین‌پویشی} سی‌تی شکمی کم بررسی شده است. آن‌چه هم انجام شده محدود است: تنها کار پارک، فیشمن و یول~\cite{park2021multiphase} مستقیماً هم‌ترازسازی پویش شریانی به وریدی را هدف گرفته است.

رویکرد آن‌ها ترکیبی از دو ایده‌ی بالاست: تغییر شکل \textbf{وابسته به جهت}، که به حرکت لغزشی در مرز اندام‌ها اجازه می‌دهد (در خط پاپیه‌ژ)، به همراه خودشباهت محلی برای مقاوم‌سازی در برابر تفاوت تباین (در خط \lr{MIND}). ارزیابی روی داده‌ی سی‌تی نازک‌برش بیماران آدنوکارسینوم پانکراس و اهداکنندگان کلیه انجام شد و کانتور دستی اندام‌ها با میدان جابه‌جایی منتقل و با کانتور دستی تصویر مقصد مقایسه شد.

نتایج، ارزشمندترین نقطه‌ی مرجع کمّی موجود برای این مسئله‌اند. روش پیشنهادی آن‌ها به \lr{Dice} برابر \عدد{۰٫۸۹۰} برای پانکراس و \عدد{۰٫۹۸۵} برای کبد رسید، با فاصله‌ی سطحی \عدد{۰٫۸۵} و \عدد{۰٫۷۸} میلی‌متر، در کمتر از یک دقیقه محاسبه. در همان ارزیابی، \lr{SyN} به \عدد{۰٫۸۷۳} و \عدد{۰٫۹۵۳}، دیمانز دیفئومورفیک به \عدد{۰٫۸۷۲} و \عدد{۰٫۹۳۵}، و \lr{MIND} به \عدد{۰٫۸۶۲} و \عدد{۰٫۹۳۹} رسیدند.

تفاوت میان کبد (بالای \عدد{۰٫۹۳} برای همه‌ی روش‌ها) و پانکراس (زیر \عدد{۰٫۹۰} برای همه) نشان می‌دهد دشواری هم‌ترازسازی بین‌پویشی، مانند دشواری بخش‌بندی در بخش~\ref{sec:seg-gap}، روی اندام‌های کوچک و کم‌تباین متمرکز است. افزون بر آن، این‌که یک روش مبتنی‌بر خودشباهت محلی از \lr{SyN} مبتنی‌بر همبستگی متقابل پیشی می‌گیرد، تأیید تجربی همان استدلالی است که در بالا مطرح شد: در حضور تغییر ناهمگن تباین، معیارهای ساختاری بر معیارهای شدت‌محور برتری دارند.

\lr{SimpleITK}~\cite{lowekamp2013} و \lr{elastix}~\cite{klein2010} پرکاربردترین زیرساخت‌های منبع‌باز این حوزه‌اند. یک ویژگی این کتابخانه‌ها از منظر مسئله‌ی این بخش اهمیت دارد: معماری پیمانه‌ای آن‌ها اجازه می‌دهد تبدیل، معیار شباهت و بهینه‌ساز به‌طور مستقل تعریف شوند، و (تعیین‌کننده‌تر) از \textbf{نقاب تصویر ثابت} در تعریف معیار پشتیبانی می‌کنند. یعنی می‌توان به کتابخانه گفت معیار شباهت را تنها روی زیرمجموعه‌ای از حجم‌دانه‌ها محاسبه کند. یعنی محدودکردن دامنه‌ی محاسبه‌ی معیار، از نظر پیاده‌سازی مانعی ندارد.

\subsection*{جمع‌بندی و جای خالی}


سه خانواده‌ی بالا، هرچه در جزئیات پیاده‌سازی متفاوت باشند، در نهایت به سه انتخاب متفاوت برای ساخت معیار تراز فرو می‌کاهند. شکل~\ref{fig:reg-archetypes} یک نماینده‌ی شاخص از هر خانواده را کنار هم می‌گذارد: نماینده‌ی خانواده‌ی نخست ترکیب تغییر شکل آزاد بی‌اسپلاین~\cite{rueckert1999} با اطلاعات متقابل مشتق‌پذیر~\cite{mattes2001} است، همان ترکیبی که هسته‌ی \lr{elastix}~\cite{klein2010} و \lr{SimpleITK}~\cite{lowekamp2013} است؛ نماینده‌ی خانواده‌ی دوم \lr{MIND}~\cite{heinrich2012} و بسط آن \lr{DEEDS}~\cite{heinrich2013} است؛ و نماینده‌ی خانواده‌ی سوم \lr{VoxelMorph}~\cite{balakrishnan2019} و جانشین ترادیس‌سازی آن \lr{TransMorph}~\cite{chen2022} است. این سه، الگوهایی را تثبیت کرده‌اند که بقیه‌ی روش‌های همان خانواده صرفاً پالایش می‌کنند: \lr{LapIRN}، \lr{HyperMorph} و شبکه‌های آبشاری همگی روی همان اسکلت پانل~(پ) بنا شده‌اند و تنها یکی از سه محدودیت \lr{VoxelMorph} را هدف می‌گیرند؛ \lr{UniGradICON} نیز همان اسکلت را با داده‌ی به‌مراتب بزرگ‌تر و یک جمله‌ی منظم‌ساز اضافی آموزش می‌دهد. به همین دلیل این سه به‌عنوان نماینده‌ی هر خانواده انتخاب شدند.

جدول~\ref{tab:reg-families} هسته‌ی معماری، توانایی اصلی و کاستی اصلی هر سه خانواده را کنار هم می‌گذارد.

\begin{table}[H]
\centering
\footnotesize
\caption{سه خانواده‌ی روش‌شناختی هم‌ترازسازی تصویر.}
\label{tab:reg-families}
\begin{tabular}{|p{2.6cm}|p{3.7cm}|p{3.4cm}|p{3.4cm}|}
\hline
\textbf{خانواده} & \textbf{هسته‌ی معماری} & \textbf{توانایی اصلی} & \textbf{کاستی اصلی} \\
\hline
معیارهای شدت‌محور (اطلاعات متقابل/همبستگی متقابل) & تابع هزینه از توزیع مشترک شدت روشنایی روی کل تصویر ساخته می‌شود & بدون نیاز به آموزش؛ پایه‌ی نرم‌افزارهای استاندارد (\lr{elastix}/\lr{ANTs}) & فرض پایداری رابطه‌ی شدت در کل میدان دید را تزریق ماده‌ی حاجب نقض می‌کند \\
\hline
توصیف‌گرهای ساختاری (\lr{MIND}/\lr{DEEDS}) & مقایسه‌ی الگوی محلی خودشباهت به‌جای مقدار خام شدت & مقاوم در برابر تغییر تباین و وجهیت؛ دقیق‌ترین گروه در محک‌های شکمی & بهینه‌سازی گسسته پرهزینه؛ تغییرشکل غیرصلب می‌تواند تاخوردگی تولید کند \\
\hline
روش‌های یادگیری‌محور (\lr{VoxelMorph} و مشتقات آن) & شبکه‌ای که میدان جابه‌جایی را در یک گذر پیش‌رو پیش‌بینی می‌کند & استنتاج در چند ثانیه؛ قابل تعمیم با آموزش روی داده‌ی بزرگ & تقریباً همگی روی محک‌های مغزی توسعه یافته‌اند و انتقال به شکم افت محسوس دارد \\
\hline
\end{tabular}
\end{table}

هر سه خانواده در یک چیز مشترک‌اند: در هر سه، معیار تراز روی \emph{کل} میدان دید تعریف می‌شود؛ چه به‌صورت جمع روی همه‌ی حجم‌دانه‌ها، چه به‌صورت جمع روی همه‌ی بردارهای توصیف‌گر، و چه به‌صورت زیانی که روی کل تصویر تغییرشکل‌یافته محاسبه می‌شود.

% ---------------------------------------------------------------
%  شکل: سه سرگروه روش‌شناختی هم‌ترازسازی تصویر
%  چیدمان با مختصات صریح (نه positioning نسبی) تا متن‌ها روی هم نیفتند.
%  نیازمند: tikz + arrows.meta
%  استفاده:  % ---------------------------------------------------------------
%  شکل: سه سرگروه روش‌شناختی هم‌ترازسازی تصویر
%  چیدمان با مختصات صریح (نه positioning نسبی) تا متن‌ها روی هم نیفتند.
%  نیازمند: tikz + arrows.meta
%  استفاده:  % ---------------------------------------------------------------
%  شکل: سه سرگروه روش‌شناختی هم‌ترازسازی تصویر
%  چیدمان با مختصات صریح (نه positioning نسبی) تا متن‌ها روی هم نیفتند.
%  نیازمند: tikz + arrows.meta
%  استفاده:  \input{figs/arch_registration}
% ---------------------------------------------------------------
\begingroup
\providecommand{\rsubtext}[1]{{\color{xtire}\tiny #1}}

\tikzset{
  rb/.style   ={draw=xnokhodi, line width=.5pt, fill=white, rounded corners=1.5pt,
                text width=2.55cm, minimum height=1.05cm, align=center,
                inner xsep=1.2mm, inner ysep=1.2mm, font=\scriptsize},
  rnet/.style ={rb, fill=xabikam, draw=xabi},
  rloss/.style={rb, fill=xkamrang},
  rio/.style  ={rb, draw=xkhakestari},
  rflow/.style={->, draw=xtire, line width=.6pt,
                >={Stealth[length=1.7mm,width=1.4mm]}},
  rback/.style={->, draw=xghermez, line width=.5pt, dash pattern=on 1.8pt off 1.4pt,
                >={Stealth[length=1.7mm,width=1.4mm]}},
  rlab/.style ={font=\tiny, color=xkhakestari, inner sep=1.2pt, fill=white},
}

\newcommand{\Rhead}[2]{%
  {\small\textbf{#1}}\par\nobreak\vspace{0.4mm}%
  {\scriptsize\color{xtire}#2}\par\nobreak\vspace{2.2mm}}

\begin{figure}[H]
\centering

% ═════════════════ (الف) کلاسیک شدت‌محور ═════════════════
\Rhead{(الف) خانواده‌ی کلاسیک شدت‌محور \,---\, \lr{elastix} / \lr{SimpleITK}}
      {بدون مرحله‌ی آموزش؛ برای هر جفت تصویر یک بهینه‌سازی تکراری جداگانه اجرا می‌شود}
\begin{tikzpicture}
  \node[rio]   (mov)  at (     0, 0)    {\beginR تصویر متحرک\endR};
  \node[rnet]  (tr)   at (  -3.5, 0)    {\beginR تبدیل $T_\theta$\endR\\[1pt]\rsubtext{\beginR شبکه‌ی نقاط کنترل بی‌اسپلاین\endR}};
  \node[rio]   (warp) at (    -7, 0)    {\beginR تصویر تغییرشکل‌یافته\endR};
  \node[rloss] (sim)  at ( -10.5, 0)    {\beginR معیار شباهت\endR\\[1pt]\rsubtext{\beginR اطلاعات متقابل\endR}};
  \node[rio]   (fix)  at ( -10.5, 1.8)  {\beginR تصویر ثابت\endR};
  \node[rnet]  (opt)  at (    -7,-1.8)  {\beginR بهینه‌ساز گرادیان‌محور\endR};

  \draw[rflow] (mov.west)  -- (tr.east);
  \draw[rflow] (tr.west)   -- (warp.east);
  \draw[rflow] (warp.west) -- (sim.east);
  \draw[rflow] (fix.south) -- (sim.north);
  \draw[rflow] (sim.south) |- (opt.west);
  \draw[rback] (opt.east)  -| node[rlab, pos=0.35] {\beginR به‌روزرسانی $\theta$ و تکرار\endR} (tr.south);
\end{tikzpicture}

\par\vspace{7mm}

% ═════════════════ (ب) توصیف‌گر ساختاری ═════════════════
\Rhead{(ب) توصیف‌گرهای ساختاری \,---\, \lr{MIND} / \lr{DEEDS}}
      {به‌جای \emph{مقدار} شدت روشنایی، \emph{الگوی محلی} آن مقایسه می‌شود}
\begin{tikzpicture}
  \node[rio]   (mv)  at (    0, 1.0)  {\beginR تصویر متحرک\endR};
  \node[rio]   (fx)  at (    0,-1.0)  {\beginR تصویر ثابت\endR};
  \node[rnet]  (d1)  at ( -3.5, 1.0)  {\beginR توصیف‌گر \lr{MIND}\endR\\[1pt]\rsubtext{\beginR خودشباهت وصله‌ها\endR}};
  \node[rnet]  (d2)  at ( -3.5,-1.0)  {\beginR توصیف‌گر \lr{MIND}\endR};
  \node[rloss] (ssd) at (   -7, 0)    {\beginR \lr{SSD} روی بردار توصیف‌گر\endR};
  \node[rnet]  (mrf) at (-10.5, 0)    {\beginR بهینه‌سازی گسسته‌ی \lr{MRF}\endR\\[1pt]\rsubtext{\beginR انتقال پیام روی گراف\endR}};
  \node[rio]   (phi) at (-10.5,-2.0)  {\beginR میدان جابه‌جایی\endR};

  \draw[rflow] (mv.west) -- (d1.east);
  \draw[rflow] (fx.west) -- (d2.east);
  \draw[rflow] (d1.west) -- ++(-0.55,0) |- (ssd.east);
  \draw[rflow] (d2.west) -- ++(-0.55,0) |- (ssd.east);
  \draw[rflow] (ssd.west)  -- (mrf.east);
  \draw[rflow] (mrf.south) -- (phi.north);
\end{tikzpicture}

\par\vspace{7mm}

% ═════════════════ (پ) یادگیری‌محور ═════════════════
\Rhead{(پ) خانواده‌ی یادگیری‌محور \,---\, \lr{VoxelMorph} / \lr{TransMorph}}
      {آموزش یک‌بار روی کل مجموعه؛ استنتاج در یک گذر پیش‌رو}
\begin{tikzpicture}
  \node[rio]   (in)   at (     0, 0)    {\beginR متحرک $\oplus$ ثابت\endR\\[1pt]\rsubtext{\beginR الحاق در بعد کانال\endR}};
  \node[rnet]  (enc)  at (  -3.5, 0)    {\beginR رمزگذار--رمزگشا\endR\\[1pt]\rsubtext{\beginR \lr{U-Net} یا \lr{Swin}\endR}};
  \node[rio]   (phi3) at (    -7, 0)    {\beginR میدان جابه‌جایی $\phi$\endR};
  \node[rnet]  (stn)  at ( -10.5, 0)    {\beginR تبدیل‌گر فضایی\endR\\[1pt]\rsubtext{مشتق‌پذیر}};
  \node[rloss] (ls)   at (    -7,-1.9)  {\beginR زیان: شباهت تصویر $+$ هموارسازی $\phi$\endR};

  \draw[rflow] (in.west)   -- (enc.east);
  \draw[rflow] (enc.west)  -- (phi3.east);
  \draw[rflow] (phi3.west) -- (stn.east);
  \draw[rflow] (stn.south) |- (ls.west);
  \draw[rback] (ls.east)   -| node[rlab, pos=0.35] {\beginR پس‌انتشار (فقط در آموزش)\endR} (enc.south);
\end{tikzpicture}

\caption{سه سرگروه روش‌شناختی هم‌ترازسازی تصویر: (الف) معیارهای شدت‌محور، (ب) توصیف‌گرهای ساختاری، (پ) روش‌های یادگیری‌محور.}
\label{fig:reg-archetypes}
\end{figure}
\endgroup

% ---------------------------------------------------------------
\begingroup
\providecommand{\rsubtext}[1]{{\color{xtire}\tiny #1}}

\tikzset{
  rb/.style   ={draw=xnokhodi, line width=.5pt, fill=white, rounded corners=1.5pt,
                text width=2.55cm, minimum height=1.05cm, align=center,
                inner xsep=1.2mm, inner ysep=1.2mm, font=\scriptsize},
  rnet/.style ={rb, fill=xabikam, draw=xabi},
  rloss/.style={rb, fill=xkamrang},
  rio/.style  ={rb, draw=xkhakestari},
  rflow/.style={->, draw=xtire, line width=.6pt,
                >={Stealth[length=1.7mm,width=1.4mm]}},
  rback/.style={->, draw=xghermez, line width=.5pt, dash pattern=on 1.8pt off 1.4pt,
                >={Stealth[length=1.7mm,width=1.4mm]}},
  rlab/.style ={font=\tiny, color=xkhakestari, inner sep=1.2pt, fill=white},
}

\newcommand{\Rhead}[2]{%
  {\small\textbf{#1}}\par\nobreak\vspace{0.4mm}%
  {\scriptsize\color{xtire}#2}\par\nobreak\vspace{2.2mm}}

\begin{figure}[H]
\centering

% ═════════════════ (الف) کلاسیک شدت‌محور ═════════════════
\Rhead{(الف) خانواده‌ی کلاسیک شدت‌محور \,---\, \lr{elastix} / \lr{SimpleITK}}
      {بدون مرحله‌ی آموزش؛ برای هر جفت تصویر یک بهینه‌سازی تکراری جداگانه اجرا می‌شود}
\begin{tikzpicture}
  \node[rio]   (mov)  at (     0, 0)    {\beginR تصویر متحرک\endR};
  \node[rnet]  (tr)   at (  -3.5, 0)    {\beginR تبدیل $T_\theta$\endR\\[1pt]\rsubtext{\beginR شبکه‌ی نقاط کنترل بی‌اسپلاین\endR}};
  \node[rio]   (warp) at (    -7, 0)    {\beginR تصویر تغییرشکل‌یافته\endR};
  \node[rloss] (sim)  at ( -10.5, 0)    {\beginR معیار شباهت\endR\\[1pt]\rsubtext{\beginR اطلاعات متقابل\endR}};
  \node[rio]   (fix)  at ( -10.5, 1.8)  {\beginR تصویر ثابت\endR};
  \node[rnet]  (opt)  at (    -7,-1.8)  {\beginR بهینه‌ساز گرادیان‌محور\endR};

  \draw[rflow] (mov.west)  -- (tr.east);
  \draw[rflow] (tr.west)   -- (warp.east);
  \draw[rflow] (warp.west) -- (sim.east);
  \draw[rflow] (fix.south) -- (sim.north);
  \draw[rflow] (sim.south) |- (opt.west);
  \draw[rback] (opt.east)  -| node[rlab, pos=0.35] {\beginR به‌روزرسانی $\theta$ و تکرار\endR} (tr.south);
\end{tikzpicture}

\par\vspace{7mm}

% ═════════════════ (ب) توصیف‌گر ساختاری ═════════════════
\Rhead{(ب) توصیف‌گرهای ساختاری \,---\, \lr{MIND} / \lr{DEEDS}}
      {به‌جای \emph{مقدار} شدت روشنایی، \emph{الگوی محلی} آن مقایسه می‌شود}
\begin{tikzpicture}
  \node[rio]   (mv)  at (    0, 1.0)  {\beginR تصویر متحرک\endR};
  \node[rio]   (fx)  at (    0,-1.0)  {\beginR تصویر ثابت\endR};
  \node[rnet]  (d1)  at ( -3.5, 1.0)  {\beginR توصیف‌گر \lr{MIND}\endR\\[1pt]\rsubtext{\beginR خودشباهت وصله‌ها\endR}};
  \node[rnet]  (d2)  at ( -3.5,-1.0)  {\beginR توصیف‌گر \lr{MIND}\endR};
  \node[rloss] (ssd) at (   -7, 0)    {\beginR \lr{SSD} روی بردار توصیف‌گر\endR};
  \node[rnet]  (mrf) at (-10.5, 0)    {\beginR بهینه‌سازی گسسته‌ی \lr{MRF}\endR\\[1pt]\rsubtext{\beginR انتقال پیام روی گراف\endR}};
  \node[rio]   (phi) at (-10.5,-2.0)  {\beginR میدان جابه‌جایی\endR};

  \draw[rflow] (mv.west) -- (d1.east);
  \draw[rflow] (fx.west) -- (d2.east);
  \draw[rflow] (d1.west) -- ++(-0.55,0) |- (ssd.east);
  \draw[rflow] (d2.west) -- ++(-0.55,0) |- (ssd.east);
  \draw[rflow] (ssd.west)  -- (mrf.east);
  \draw[rflow] (mrf.south) -- (phi.north);
\end{tikzpicture}

\par\vspace{7mm}

% ═════════════════ (پ) یادگیری‌محور ═════════════════
\Rhead{(پ) خانواده‌ی یادگیری‌محور \,---\, \lr{VoxelMorph} / \lr{TransMorph}}
      {آموزش یک‌بار روی کل مجموعه؛ استنتاج در یک گذر پیش‌رو}
\begin{tikzpicture}
  \node[rio]   (in)   at (     0, 0)    {\beginR متحرک $\oplus$ ثابت\endR\\[1pt]\rsubtext{\beginR الحاق در بعد کانال\endR}};
  \node[rnet]  (enc)  at (  -3.5, 0)    {\beginR رمزگذار--رمزگشا\endR\\[1pt]\rsubtext{\beginR \lr{U-Net} یا \lr{Swin}\endR}};
  \node[rio]   (phi3) at (    -7, 0)    {\beginR میدان جابه‌جایی $\phi$\endR};
  \node[rnet]  (stn)  at ( -10.5, 0)    {\beginR تبدیل‌گر فضایی\endR\\[1pt]\rsubtext{مشتق‌پذیر}};
  \node[rloss] (ls)   at (    -7,-1.9)  {\beginR زیان: شباهت تصویر $+$ هموارسازی $\phi$\endR};

  \draw[rflow] (in.west)   -- (enc.east);
  \draw[rflow] (enc.west)  -- (phi3.east);
  \draw[rflow] (phi3.west) -- (stn.east);
  \draw[rflow] (stn.south) |- (ls.west);
  \draw[rback] (ls.east)   -| node[rlab, pos=0.35] {\beginR پس‌انتشار (فقط در آموزش)\endR} (enc.south);
\end{tikzpicture}

\caption{سه سرگروه روش‌شناختی هم‌ترازسازی تصویر: (الف) معیارهای شدت‌محور، (ب) توصیف‌گرهای ساختاری، (پ) روش‌های یادگیری‌محور.}
\label{fig:reg-archetypes}
\end{figure}
\endgroup

% ---------------------------------------------------------------
\begingroup
\providecommand{\rsubtext}[1]{{\color{xtire}\tiny #1}}

\tikzset{
  rb/.style   ={draw=xnokhodi, line width=.5pt, fill=white, rounded corners=1.5pt,
                text width=2.55cm, minimum height=1.05cm, align=center,
                inner xsep=1.2mm, inner ysep=1.2mm, font=\scriptsize},
  rnet/.style ={rb, fill=xabikam, draw=xabi},
  rloss/.style={rb, fill=xkamrang},
  rio/.style  ={rb, draw=xkhakestari},
  rflow/.style={->, draw=xtire, line width=.6pt,
                >={Stealth[length=1.7mm,width=1.4mm]}},
  rback/.style={->, draw=xghermez, line width=.5pt, dash pattern=on 1.8pt off 1.4pt,
                >={Stealth[length=1.7mm,width=1.4mm]}},
  rlab/.style ={font=\tiny, color=xkhakestari, inner sep=1.2pt, fill=white},
}

\newcommand{\Rhead}[2]{%
  {\small\textbf{#1}}\par\nobreak\vspace{0.4mm}%
  {\scriptsize\color{xtire}#2}\par\nobreak\vspace{2.2mm}}

\begin{figure}[H]
\centering

% ═════════════════ (الف) کلاسیک شدت‌محور ═════════════════
\Rhead{(الف) خانواده‌ی کلاسیک شدت‌محور \,---\, \lr{elastix} / \lr{SimpleITK}}
      {بدون مرحله‌ی آموزش؛ برای هر جفت تصویر یک بهینه‌سازی تکراری جداگانه اجرا می‌شود}
\begin{tikzpicture}
  \node[rio]   (mov)  at (     0, 0)    {\beginR تصویر متحرک\endR};
  \node[rnet]  (tr)   at (  -3.5, 0)    {\beginR تبدیل $T_\theta$\endR\\[1pt]\rsubtext{\beginR شبکه‌ی نقاط کنترل بی‌اسپلاین\endR}};
  \node[rio]   (warp) at (    -7, 0)    {\beginR تصویر تغییرشکل‌یافته\endR};
  \node[rloss] (sim)  at ( -10.5, 0)    {\beginR معیار شباهت\endR\\[1pt]\rsubtext{\beginR اطلاعات متقابل\endR}};
  \node[rio]   (fix)  at ( -10.5, 1.8)  {\beginR تصویر ثابت\endR};
  \node[rnet]  (opt)  at (    -7,-1.8)  {\beginR بهینه‌ساز گرادیان‌محور\endR};

  \draw[rflow] (mov.west)  -- (tr.east);
  \draw[rflow] (tr.west)   -- (warp.east);
  \draw[rflow] (warp.west) -- (sim.east);
  \draw[rflow] (fix.south) -- (sim.north);
  \draw[rflow] (sim.south) |- (opt.west);
  \draw[rback] (opt.east)  -| node[rlab, pos=0.35] {\beginR به‌روزرسانی $\theta$ و تکرار\endR} (tr.south);
\end{tikzpicture}

\par\vspace{7mm}

% ═════════════════ (ب) توصیف‌گر ساختاری ═════════════════
\Rhead{(ب) توصیف‌گرهای ساختاری \,---\, \lr{MIND} / \lr{DEEDS}}
      {به‌جای \emph{مقدار} شدت روشنایی، \emph{الگوی محلی} آن مقایسه می‌شود}
\begin{tikzpicture}
  \node[rio]   (mv)  at (    0, 1.0)  {\beginR تصویر متحرک\endR};
  \node[rio]   (fx)  at (    0,-1.0)  {\beginR تصویر ثابت\endR};
  \node[rnet]  (d1)  at ( -3.5, 1.0)  {\beginR توصیف‌گر \lr{MIND}\endR\\[1pt]\rsubtext{\beginR خودشباهت وصله‌ها\endR}};
  \node[rnet]  (d2)  at ( -3.5,-1.0)  {\beginR توصیف‌گر \lr{MIND}\endR};
  \node[rloss] (ssd) at (   -7, 0)    {\beginR \lr{SSD} روی بردار توصیف‌گر\endR};
  \node[rnet]  (mrf) at (-10.5, 0)    {\beginR بهینه‌سازی گسسته‌ی \lr{MRF}\endR\\[1pt]\rsubtext{\beginR انتقال پیام روی گراف\endR}};
  \node[rio]   (phi) at (-10.5,-2.0)  {\beginR میدان جابه‌جایی\endR};

  \draw[rflow] (mv.west) -- (d1.east);
  \draw[rflow] (fx.west) -- (d2.east);
  \draw[rflow] (d1.west) -- ++(-0.55,0) |- (ssd.east);
  \draw[rflow] (d2.west) -- ++(-0.55,0) |- (ssd.east);
  \draw[rflow] (ssd.west)  -- (mrf.east);
  \draw[rflow] (mrf.south) -- (phi.north);
\end{tikzpicture}

\par\vspace{7mm}

% ═════════════════ (پ) یادگیری‌محور ═════════════════
\Rhead{(پ) خانواده‌ی یادگیری‌محور \,---\, \lr{VoxelMorph} / \lr{TransMorph}}
      {آموزش یک‌بار روی کل مجموعه؛ استنتاج در یک گذر پیش‌رو}
\begin{tikzpicture}
  \node[rio]   (in)   at (     0, 0)    {\beginR متحرک $\oplus$ ثابت\endR\\[1pt]\rsubtext{\beginR الحاق در بعد کانال\endR}};
  \node[rnet]  (enc)  at (  -3.5, 0)    {\beginR رمزگذار--رمزگشا\endR\\[1pt]\rsubtext{\beginR \lr{U-Net} یا \lr{Swin}\endR}};
  \node[rio]   (phi3) at (    -7, 0)    {\beginR میدان جابه‌جایی $\phi$\endR};
  \node[rnet]  (stn)  at ( -10.5, 0)    {\beginR تبدیل‌گر فضایی\endR\\[1pt]\rsubtext{مشتق‌پذیر}};
  \node[rloss] (ls)   at (    -7,-1.9)  {\beginR زیان: شباهت تصویر $+$ هموارسازی $\phi$\endR};

  \draw[rflow] (in.west)   -- (enc.east);
  \draw[rflow] (enc.west)  -- (phi3.east);
  \draw[rflow] (phi3.west) -- (stn.east);
  \draw[rflow] (stn.south) |- (ls.west);
  \draw[rback] (ls.east)   -| node[rlab, pos=0.35] {\beginR پس‌انتشار (فقط در آموزش)\endR} (enc.south);
\end{tikzpicture}

\caption{سه سرگروه روش‌شناختی هم‌ترازسازی تصویر: (الف) معیارهای شدت‌محور، (ب) توصیف‌گرهای ساختاری، (پ) روش‌های یادگیری‌محور.}
\label{fig:reg-archetypes}
\end{figure}
\endgroup


معیارهای کلاسیک خانواده‌ی اطلاعات متقابل در سطح کل تصویر تعریف می‌شوند و بنابراین در برابر تغییر ناهمگن تباین شکننده‌اند. روش‌های یادگیری‌محور سریع‌ترند اما تقریباً همگی روی محک‌های مغزی توسعه یافته‌اند و (چنان‌که افت \lr{Dice} از \عدد{۰٫۷۶۲} به \عدد{۰٫۴۸۳} در \lr{UniGradICON} نشان می‌دهد) انتقالشان به شکم بدیهی نیست. توصیف‌گرهای ساختاری مانند \lr{MIND} و \lr{DEEDS} مقاوم‌ترین گزینه‌اند، و کار پارک و همکاران~\cite{park2021multiphase} تأیید می‌کند که همین خانواده در هم‌ترازسازی بین‌پویشی شکمی بهترین عملکرد را دارد.

جای خالی، از این رو، این نیست که کسی هم‌ترازسازی بین‌پویشی شکمی انجام نداده باشد؛ پارک و همکاران~\cite{park2021multiphase} انجام داده‌اند. جای خالی باریک‌تر است: همه‌ی روش‌های موجود، چه آن‌ها که آناتومی را برای \emph{وزن‌دهی} به کار می‌گیرند (\cite{papiez2014,devos2019,oktay2017}) و چه آن‌ها که معیار را \emph{مقاوم‌تر} می‌سازند (\cite{heinrich2012,heinrich2013,park2021multiphase})، معیار تراز را همچنان روی \emph{کل} میدان دید محاسبه می‌کنند؛ هیچ‌یک این دامنه را از پیش، بر مبنای دانش تشریحی، محدود نمی‌کند.

افزون بر این، نکته‌ی مشترک دیگری این ادبیات را به مسئله‌ی گسترده‌تر این فصل گره می‌زند. در سراسر ادبیات هم‌نهادی و هم‌ترازسازی، کیفیت هم‌ترازی جفت‌های آموزشی هیچ‌گاه \emph{اندازه‌گیری و گزارش} نشده است. مقالات یا می‌گویند داده را هم‌تراز کرده‌اند (اغلب با ذکر نام یک ابزار و بدون هیچ عددی) یا اصلاً به آن اشاره نمی‌کنند. پیامد این سکوت آن است که خواننده نمی‌داند هر مدل با داده‌ای به چه کیفیتی آموزش دیده است؛ و چون هیچ عدد مشترکی وجود ندارد، مقایسه‌ی ارقام هم‌نهادی میان مقالات، سهم \emph{مدل} را از سهم \emph{کیفیت داده} تفکیک نمی‌کند. به بیان دیگر، در این ادبیات هم‌ترازسازی همواره \emph{هدف} است و هیچ‌گاه به‌عنوان \emph{معیار پذیرش داده} برای یک وظیفه‌ی پایین‌دستی به کار نرفته است؛ کاربردی که به یک آستانه‌ی کمّی نیاز دارد و چنین آستانه‌ای در ادبیات موجود تعریف نشده است.


\section{تولید تصاویر سی‌تی با ماده‌ی حاجب از تصاویر سی‌تی بدون تزریق ماده‌ی حاجب}
\label{sec:synthesis}

این بخش هسته‌ی فصل است. روش‌های این حوزه با معماری‌ها، مجموعه‌داده‌ها و معیارهای بسیار متفاوت گزارش شده‌اند و در این بخش بر اساس راهبرد نظارتی خود به پنج خانواده تقسیم می‌شوند. پس از آن معماری‌های عمومی‌ای که مبنای متعارف این ادبیات شده‌اند مرور می‌شوند و بخش با سه سرگروه معماری و جدول کامل روش‌ها پایان می‌یابد.


\subsubsection*{خانواده‌ی نخست: \lr{GAN}های جفت‌شده}

نقطه‌ی عزیمت این ادبیات، معماری \lr{pix2pix} است: یک \lr{GAN} شرطی که تصویر ورودی را به تصویر خروجی نگاشت می‌کند و با ترکیبی از زیان تخاصمی و زیان بازسازی تصویردانه‌ای آموزش می‌بیند. مطالعات این خانواده بیشتر در یک نکته با هم فرق دارند: روشی که مولد را وادار می‌کنند به \emph{نواحی مرتبط با ماده‌ی حاجب} توجه کند، نه به کل تصویر به‌طور یکنواخت.

مجموعه‌ای از سه کار پیاپی از یانگ و همکاران، این تکامل را به‌روشنی نشان می‌دهد؛ هر سه روی یک مجموعه‌داده‌ی درون‌سازمانی مشترک از \عدد{۸۶} بیمار با چهار نوع ضایعه‌ی کبدی (کارسینوم سلول کبدی، هیپرپلازی ندولار کانونی، متاستاز و کیست) ارزیابی شده‌اند که برای هر بیمار سه پویش بدون‌تزریق، شریانی و وریدی دارد.

در گام نخست~\cite{Yan22}، یک \textbf{\lr{cGAN} بهبودیافته} معرفی شد که دو ضعف مشخص معماری پایه را هدف می‌گیرد: میدان دریافت محدود رمزگذار پیچشی، و رقیق‌شدن اطلاعات معنایی در مسیر رمزگشایی. راه‌حل، افزودن یک \textbf{پیمانه تجمیع هرمی}\footnote{\lr{Pyramid pooling module}} به انتهای رمزگذار و مجموعه‌ای از مسیرهای جریان اطلاعات است که خروجی این پیمانه را مستقیماً به هر واحد رمزگشا می‌رساند. رمزگذار از \lr{ResNet} از پیش‌آموزش‌دیده روی \lr{ImageNet} استفاده می‌کند و شبکه با \lr{WGAN-GP} به‌همراه زیان بازسازی \lr{L1} (با وزن \عدد{۲۵} در برابر وزن \عدد{۱} برای جمله‌ی تخاصمی) آموزش می‌بیند. نتیجه، \lr{PSNR} برابر \عدد{۳۳٫۱۸} و \lr{SSIM} برابر \عدد{۰٫۸۱۹} برای پویش شریانی و \عدد{۳۳٫۰۲} و \عدد{۰٫۸۱۹} برای پویش وریدی بود، در برابر \عدد{۳۲٫۸۶} و \عدد{۰٫۷۹۷} برای \lr{pix2pix} پایه.

گام دوم~\cite{Yan23}، با نام \textbf{\lr{SAG-GAN}}، مسئله‌ی توجه را صریح کرد. یک \textbf{پیمانه توجه فضایی} سبک، نقشه‌ی امتیاز فضایی تولید می‌کند که مشخص می‌کند کدام نواحی تصویر با وظیفه مرتبط‌اند؛ تصویر نهایی، ترکیب وزن‌دار تصویر تبدیل‌شده و تصویر ورودی بر اساس همین نقشه است، به‌طوری‌که نواحی نامرتبط دست‌نخورده باقی می‌مانند. نکته‌ی هوشمندانه آن است که این نقشه بدون نیاز به نقاب دستی و به‌صورت \emph{بدون‌نظارت} آموخته می‌شود، از طریق یک \textbf{زیان توجه} که سازگاری نقشه‌ی امتیاز میان تصویر ورودی و تصویر هم‌نهاده را تشویق می‌کند. تابع هدف نهایی ترکیبی سه‌جزئی است (تخاصمی با وزن \عدد{۲}، بازسازی با وزن \عدد{۲۵}، توجه با وزن \عدد{۱})، و \lr{SSIM} پویش شریانی را به \عدد{۰٫۸۳۰} و ضریب همبستگی پیرسون ناحیه‌ی تومور را به \عدد{۰٫۹۶۶} رساند.

گام سوم~\cite{Yan24c}، \textbf{\lr{SGCDD-GAN}}، دو تحول افزود. نخست، \textbf{رمزگشایی دوگانه‌ی متقاطع}: یک شاخه تصویر تبدیل‌شده و شاخه‌ی دیگر نقشه‌ی توجه را تولید می‌کند، اما این دو شاخه به‌صورت متقاطع به هم متصل‌اند تا اطلاعات میانشان جریان یابد. دوم، \textbf{یادگیری چندوظیفه‌ای}: یک شبکه‌ی بخش‌بندی تومور به چارچوب افزوده می‌شود تا مولد را وادار کند به ناحیه‌ی ضایعه توجه بیشتری کند. ارزیابی این کار، برخلاف دو کار پیشین، از معیارهای سطح‌تصویردانه فراتر رفت: تصاویر هم‌نهاده به یک دسته‌بند \lr{ResNet-50} سه‌پویشی داده شد و به دقت \عدد{۰٫۸۲۷} در دسته‌بندی نوع ضایعه رسید (در برابر \عدد{۰٫۶۳۵} با پویش بدون‌تزریق تنها). در یک آزمایش مقدماتی، دو رادیولوژیست با تصاویر هم‌نهاده به دقت تشخیصی \عدد{۰٫۹۴۱} رسیدند، در برابر \عدد{۱٫۰۰} با تصویر واقعی.

\textbf{\lr{MAN-GAN}}~\cite{Zha25} مسیر مکملی را دنبال کرد: به‌جای آموختن توجه، نقاب اندام کبد را \emph{مستقیماً} از طریق \textbf{بهنجارسازی سازگار با نقاب}\footnote{\lr{Mask-adaptive normalisation}} وارد مولد کرد؛ سازوکاری که پراسنجه\footnote{\lr{Parameter}}های بهنجارسازی لایه را بسته به این‌که حجم‌دانه درون اندام هدف است یا بیرون آن، متفاوت تنظیم می‌کند. این مطالعه گسترده‌ترین ارزیابی این خانواده را دارد: \عدد{۳۷۴} بیمار برای آموزش، \عدد{۱۴۴} بیمار با یک بیماری کبدی مشخص به‌عنوان آزمون داخلی، و \عدد{۸۳} بیمار برای اعتبارسنجی خارجی. تصویر بدون‌تزریق خام، پیش از هم‌نهادی، \lr{PSNR} برابر \عدد{۲۶٫۰۲} و \lr{SSIM} برابر \عدد{۰٫۷۴} دارد؛ یعنی کل بهبود مدل حدود \عدد{۱٫۷} دسی‌بل روی یک مبنای از پیش بالاست.

\textbf{\lr{MAFormer}}~\cite{Che24} همین ایده را با یک ترادیس‌ساز\footnote{\lr{Transformer}} بازسازی کرد و دانش بالینی متفاوتی را وارد مسئله ساخت: رادیولوژیست‌ها برای برجسته‌کردن نواحی تقویت‌شده، تصویر بدون‌تزریق را از تصویر با ماده‌ی حاجب \emph{تفریق} می‌کنند. \lr{MAFormer} یک تخمین‌گر نقاب دارد که همین تصویر تفریق را از روی تصویر بدون‌تزریق پیش‌بینی می‌کند، و آن را از طریق یک لایه‌ی بهنجارسازی مبتنی‌بر ترادیس‌ساز به شبکه تزریق می‌کند. یک \textbf{زیان تغییرناپذیری ساختار} نیز تضمین می‌کند آناتومی در برابر تغییر رنگ‌آمیزی ناشی از ماده‌ی حاجب حفظ شود. ارقام اصلی \lr{MAFormer} مربوط به مجموعه‌داده‌ی خصوصی \lr{MPECT-data} است؛ همین مطالعه روی مجموعه‌داده‌ی عمومی \lr{VinDr-Multiphase}~\cite{vindr} (همان مجموعه‌داده‌ی این پایان‌نامه) نیز ارزیابی شده و به \lr{PSNR} \عدد{۲۱٫۷۱} (شریانی) و \عدد{۲۱٫۹۹} (وریدی) و \lr{SSIM} \عدد{۰٫۷۴۹} (هر دو پویش) رسیده است؛ افتی محسوس نسبت به ارقام \lr{MPECT-data} که دوباره ناهمگنی دامنه‌ی معیارها میان مجموعه‌داده‌ها را نشان می‌دهد.

\textbf{\lr{AMNet}}~\cite{Zhu25c} تازه‌ترین نماینده‌ی این خانواده است و بر دو چالش تمرکز دارد: تباین پایین و حساسیت پایین به ماده‌ی حاجب. راه‌حل آن چهار مؤلفه دارد: یک \textbf{ترادیس‌ساز تجمیع‌شده--تعاملی} با تجمیع توکن چندمقیاسی و تعامل میان‌توکنی؛ \textbf{یادگیری متضاد هم‌ترازشده‌ی چندگرانی} که ویژگی‌های درون‌حوزه را فشرده و میان‌حوزه را جداشدنی می‌کند؛ یک \textbf{لایه‌ی سازگار آگاه از تباین}؛ و یک \textbf{متمایزکننده‌ی دوجریانه}. تابع هدف ترکیبی از زیان تصویردانه‌ای \lr{L1} (وزن \عدد{۵۰})، زیان ادراکی مبتنی‌بر \lr{VGG-19} (وزن \عدد{۱۰})، زیان تخاصمی و زیان متضاد است. داده‌ی آن \عدد{۹۲} حجم سی‌تی از کلینیک کلیولند با دو پویش بدون‌تزریق و وریدی است که \emph{پیش از آموزش هم‌تراز شده‌اند}.

\textbf{\lr{pix2pix} سه‌بعدی}~\cite{Cho21} اهمیت ویژه‌ای دارد، چون یکی از معدود کارهای این حوزه با کد باز است. چوی و همکاران تباین را از سی‌تی قفسه‌سینه‌ی بدون‌تزریق هم‌نهادی کردند و (برخلاف اغلب مطالعات این خانواده) ارزیابی را در دو سطح جداگانه انجام دادند: \عدد{۲۵} بیمار برای سنجش فنی و \عدد{۱۲} بیمار برای سنجش کاربردپذیری بالینی. در سطح فنی، تصویر هم‌نهاده در مقایسه با تصویر بدون‌تزریق خام بهبود معناداری در همه‌ی معیارها نشان داد: \lr{MAE} از \عدد{۴۸٫۷۴} به \عدد{۴۱٫۷۲}، \lr{PSNR} از \عدد{۱۵٫۹۷} به \عدد{۱۷٫۴۴}، \lr{MS-SSIM} از \عدد{۰٫۸۱} به \عدد{۰٫۸۴} و \lr{LPIPS} از \عدد{۰٫۱۵} به \عدد{۰٫۱۴}. اما سطح بالینی، یافته‌ی به‌مراتب چشمگیرتری داد: نسبت تباین‌به‌نوفه غدد لنفاوی میانی‌سینه‌ای از \عدد{۰٫۷۴} به \عدد{۶٫۱۵} جهش کرد؛ بهبودی حدوداً هشت‌برابری در سنجه‌ای که مستقیماً به قابلیت تشخیص مربوط است، در حالی که \lr{PSNR} تنها حدود \عدد{۹} درصد بهتر شده بود.

این عدم‌تناسب، مهم‌ترین درس این خانواده است: \emph{معیارهای سطح‌تصویردانه، بهبود بالینی را به‌شدت کم‌برآورد می‌کنند}، چون تفاوت را روی کل تصویر میانگین می‌گیرند در حالی که کل ارزش تشخیصی در نواحی کوچکی متمرکز است.

\subsubsection*{خانواده‌ی دوم: سازگاری چرخه‌ای و آموزش بدون‌جفت}

مسئله‌ی عملی خانواده‌ی نخست آن است که به جفت‌های هم‌ترازِ هم‌بیمار نیاز دارد، همان چیزی که فصل~\ref{chap:survey} کمیابی‌اش را مستند کرد. خانواده‌ی دوم این نیاز را با \textbf{سازگاری چرخه‌ای} حذف می‌کند: دو مولد در دو جهت مخالف آموزش می‌بینند و قید اصلی این است که ترجمه‌ی رفت و برگشت باید تصویر اولیه را بازتولید کند. این قید، آموزش را با داده‌ی \emph{بدون‌جفت} ممکن می‌سازد.

\textbf{\lr{Res-CycleGAN}}~\cite{Xie21} رمزگذار--رمزگشای باقیمانده را در چارچوب \lr{CycleGAN} برای هم‌نهادی سی‌تی شکمی (بیماران کبدی، برای برنامه‌ریزی پرتودرمانی) به کار برد؛ یک استثنای روش‌شناختی در همین خانواده، چرا که برخلاف تعریف بالا، داده‌ی آموزش آن \emph{جفت‌شده} است. نویسندگان صراحتاً معماری \lr{CycleGAN} را برای مقاومت در برابر خطای باقی‌مانده‌ی هم‌ترازی انتخاب کرده‌اند، نه برای دور زدن نیاز به جفت. \textbf{\lr{CyTran}}~\cite{Ris21} لایه‌های پیچشی و توجه چندسر را در یک ترادیس‌ساز سازگار چرخه‌ای ترکیب کرد و همراه خود مجموعه‌داده‌ی \lr{Coltea-Lung-CT-100W} را منتشر کرد. شن و همکاران~\cite{Shen21} \lr{CycleGAN} و \lr{UNIT} را نه برای تولید تصویر زیبا، بلکه به‌عنوان \emph{انتقال حوزه} برای بهبود بخش‌بندی چندعضوی به کار بردند و معیارشان \lr{Dice} پایین‌دستی بود، نه \lr{PSNR}.

بهای این آزادی، از دست دادن نظارت تصویردانه‌به‌تصویردانه است: هیچ تضمینی وجود ندارد که مولد، آناتومی را حفظ کند و صرفاً بافت را «واقعی‌نما» نکند؛ خطری که در تصویربرداری پزشکی می‌تواند به \textbf{توهم}\footnote{\lr{Hallucination}} ساختارهای ناموجود بینجامد.

\subsubsection*{خانواده‌ی سوم: مدل‌های انتشار}

موج اخیر این حوزه از \lr{GAN} به \textbf{مدل‌های انتشار} گذر کرده است؛ خانواده‌ای که به‌جای آموزش تخاصمی، فرایند تدریجی حذف نوفه را می‌آموزد و به همین دلیل آموزش پایدارتر و تنوع نمونه‌ی بهتری دارد، به قیمت استنتاج کندتر.

\textbf{\lr{CFPS-Diff}}~\cite{Zhe25} کامل‌ترین نماینده‌ی مستندشده‌ی این خانواده است. دو ایده‌ی اصلی دارد: \textbf{یادگیری الگوی جریان تباین}، که به‌جای نگاشت مستقل هر پویش، پویایی گذر ماده‌ی حاجب میان پویش‌ها را مدل می‌کند؛ و \textbf{یادگیری ویژگی خاص هر پویش} به‌همراه یک زیان آنتروپی متقاطع وزن‌دار با فاصله‌ی دسته\footnote{\lr{Class}} که تمایز میان پویش‌ها را تقویت می‌کند. تابع هدف ترکیبی از \lr{MAE} و زیان موجک است. آموزش روی \عدد{۵۴۷} بیمار از یک بیمارستان و آزمون روی \عدد{۱۱۶} مورد از دو بیمارستان دیگر، به‌همراه یک مجموعه‌ی خارجی \عدد{۲۸} موردی با آسیب‌شناسی‌های نادر. نتایج روی مجموعه‌ی آزمون: میانگین \lr{MAE} برابر \عدد{۳۱٫۲۰}، \lr{SSIM} برابر \عدد{۰٫۸۲۵} و \lr{PSNR} برابر \عدد{۲۹٫۳۹}؛ و روی مجموعه‌ی خارجی به‌ترتیب \عدد{۳۰٫۷۵}، \عدد{۰٫۸۲۶} و \عدد{۲۹٫۴۲}؛ یعنی عملاً بدون افت، که برای این حوزه غیرمعمول است.

نکته‌ی روش‌شناختی مهم این مقاله آن است که پیش‌پردازش آن صراحتاً شامل \emph{هم‌ترازسازی تصویر} است: بهنجارسازی فضایی آفین با \lr{FSL} و سپس \lr{SyN} با \lr{ANTs} برای هم‌ترازکردن پویش‌های دارای ماده‌ی حاجب با پویش بدون‌تزریق. یعنی حتی جدیدترین مدل‌های انتشار نیز نمی‌توانند از پیش‌نیاز هم‌ترازی بگریزند.

سایر کارهای این خانواده؛ \lr{AdverDM}~\cite{Zhu25}، مدل کنترل تبدیل فضایی~\cite{Wan25b}، مولد فضازمانی علّی~\cite{Zhu25b} و مدل انتشار پل براونی با شرط‌دهی متنی~\cite{Li26}؛ همگی در سال‌های \عدد{۲۰۲۵} و \عدد{۲۰۲۶} منتشر شده‌اند و هنوز جزئیات کامل مجموعه‌داده و معیار ارزیابی را گزارش نکرده‌اند.

\subsubsection*{خانواده‌ی چهارم: هم‌نهادی هدایت‌شده با هم‌ترازسازی تصویر}

زیرمجموعه‌ای از این ادبیات، به‌جای فرض‌کردن هم‌ترازی، آن را بخشی از مسئله می‌داند.

\textbf{\lr{CKAP-Net}}~\cite{Zho24} کامل‌ترین نماینده‌ی این خانواده است. ایده‌ی محوری آن \textbf{تفکیک دانش تباین} است: یک شبکه‌ی جداگانه می‌آموزد چه چیزی میان دو پویش \emph{مشترک} است (آناتومی) و چه چیزی \emph{متفاوت} است (اثر ماده‌ی حاجب)، تا مولد بتواند تنها مؤلفه‌ی دوم را اضافه کند. به آن یک \textbf{ترادیس‌ساز ادراک تشریحی چندمقیاسی} و یک \textbf{زیان ادراکی مبتنی‌بر ماده‌ی حاجب} افزوده شده که (برخلاف زیان ادراکی متعارف مبتنی‌بر \lr{VGG} آموزش‌دیده روی تصاویر طبیعی دوبعدی) مخصوص همین وظیفه آموزش دیده است. داده: \عدد{۴۰۳} بیمار از مرکز اول و \عدد{۳۸} بیمار از مرکز دوم، همگی با جفت‌های بدون‌تزریق و پویش شریانی دیررس، و \عدد{۱۳} اندام شکمی که برای زیرمجموعه‌ای از بیماران به‌دست متخصصان با بیش از ده سال تجربه برچسب خورده است.

نتایج: \lr{MAE} برابر \عدد{۲۵٫۹۶}، \lr{SSIM} برابر \عدد{۰٫۸۵۵} و \lr{PSNR} برابر \عدد{۳۲٫۶۰} (در برابر \عدد{۳۸٫۱۸}، \عدد{۰٫۷۸۰} و \عدد{۲۹٫۲۸} برای مبنا). اما ارزشمندترین عدد این مقاله، معیار پایین‌دستی آن است: بخش‌بندی چندعضوی روی تصویر هم‌نهاده به \lr{Dice} برابر \عدد{۰٫۸۱۲} و فاصله‌ی هاسدورف \عدد{۵٫۴۶} میلی‌متر رسید. این، نمونه‌ی روشنی از زنجیره‌ی استدلالی «هم‌نهادی به‌عنوان پیش‌نیاز یک وظیفه‌ی پایین‌دستی» است.

ژو و همکاران~\cite{Zho23} هم‌ترازسازی تغییرشکل‌پذیر و هم‌نهادی را در یک چارچوب یادگیری چندوظیفه‌ای \emph{هم‌زمان} آموزش دادند، به‌طوری‌که خطای هم‌ترازی و خطای هم‌نهادی یکدیگر را اصلاح کنند. لی و همکاران~\cite{Li24} با پیمانه تفکیک سازگاری آناتومیک و یک زیرپیمانه هم‌ترازسازی مبتنی‌بر \lr{VoxelMorph}، محتوای تشریحی را از سبک تباین جدا کردند؛ روی داده‌ی درون‌سازمانی \عدد{۲۷۸} بیمار مبتلا به سرطان کبد به \lr{MAE} برابر \عدد{۱۶٫۳۸}، \lr{PSNR} برابر \عدد{۳۶٫۵۲} و \lr{SSIM} برابر \عدد{۰٫۸۹۲} رسیدند. باید توجه داشت که جهت هم‌نهادی در این مقاله \lr{CECT}-به-\lr{CT} است، یعنی \emph{معکوس} جهت \lr{NC}-به-\lr{CE}؛ ارزش این کار روش‌شناختی است و نه مستقیم.

اوهم و همکاران~\cite{Uhm22} با اعمال \lr{DEEDS}~\cite{heinrich2013} روی تمام پویش‌ها \emph{پیش از} هم‌نهادی، وابستگی مستقیم کیفیت هم‌نهادی به دقت هم‌ترازی پیشین را نشان دادند. این خط پژوهشی به‌روشنی تأیید می‌کند هم‌ترازسازی تصویر و هم‌نهادی، دو رویکرد رقیب نیستند بلکه دو مرحله‌ی مکمل یک زنجیره‌ی پردازشی‌اند.

\subsubsection*{خانواده‌ی پنجم: کاربردپذیری بالینی و کاهش دوز}

شماری از مطالعات، به‌جای بهینه‌کردن معیارهای تصویری، مستقیماً می‌پرسند آیا این هم‌نهادی در بالین به کار می‌آید.

هان و همکاران~\cite{Han24} بزرگ‌ترین اعتبارسنجی بالینی این حوزه را انجام دادند. آن‌ها یک \lr{U-Net} دوونیم‌بعدی را با \textbf{یادگیری ضعیف‌نظارت} و با استفاده از تصاویر بدون‌تزریق مجازی حاصل از سی‌تی دوانرژی آموزش دادند؛ راهکاری هوشمندانه برای دورزدن نیاز به جفت‌های واقعی. آموزش روی \عدد{۲٬۲۰۲} جفت و اعتبارسنجی خارجی روی \عدد{۳۹۸} بیمار از سه مرکز و چند سازنده‌ی مختلف. \lr{PSNR} برابر \عدد{۴۳٫۲۵} و \lr{SSIM} برابر \عدد{۰٫۹۲} گزارش شد، اما یافته‌ی مهم‌تر بالینی است: با افزودن تصویر هم‌نهاده به تصویر بدون‌تزریق، حساسیت سه رادیولوژیست در تشخیص ضایعه از \عدد{۰٫۷۲۰} به \عدد{۰٫۷۶۴} و سطح اطمینان تشخیصی از \عدد{۳٫۰} به \عدد{۳٫۶} افزایش یافت.

کیم و همکاران~\cite{Kim21} با یک \lr{cGAN} دومرحله‌ای، دقت تشخیصی سی‌تی مجازی با ماده‌ی حاجب را در بیماران اورژانس با درد حاد شکمی سنجیدند.

هاوبولد و همکاران~\cite{Hau21} مسئله را وارونه کردند: به‌جای حذف کامل ماده‌ی حاجب، مقدار آن را \emph{کم} کردند. آن‌ها با سی‌تی دوانرژی \عدد{۸۵} بیمار، دوز کاهش‌یافته را به‌صورت مصنوعی شبیه‌سازی کردند (با تفریق نسبی نقشه‌ی ید) و چهار \lr{cGAN} را برای بازسازی تصویر با دوز کامل آموزش دادند. مدل‌های کاهش \عدد{۵۰} درصدی به \lr{SSIM} بالای \عدد{۰٫۹۹} و \lr{PSNR} بالای \عدد{۵۱} رسیدند و مدل‌های کاهش \عدد{۸۰} درصدی به \lr{SSIM} بالای \عدد{۰٫۹۸} و \lr{PSNR} بالای \عدد{۴۸}.

اما نتیجه‌گیری بالینی آن‌ها محتاطانه‌تر از این ارقام است و درسی مهم در بر دارد: تنها شبکه‌های کاهش \عدد{۵۰} درصدی به \textbf{سازگاری آسیب‌شناختی} \عدد{۱۰۰} درصدی رسیدند؛ در سطح کاهش \عدد{۸۰} درصدی، با وجود \lr{SSIM} بالای \عدد{۰٫۹۸}، یافته‌های آسیب‌شناختی در برخی موارد تغییر کردند. رقم \عدد{۷۰} درصدی کاهش دوز که گاه به این خط پژوهشی نسبت داده می‌شود، متعلق به مطالعه‌ی \emph{حیوانی} بعدی همین گروه~\cite{Hau22} روی \عدد{۲۰} مینی‌پیگ است و نه به مطالعه‌ی انسانی حاضر؛ در آن مطالعه نسبت تباین‌به‌نوفه عروقی از \عدد{۶٫۹} به \عدد{۲۸٫۷} رسید (مقدار واقعی با دوز کامل \عدد{۳۴٫۲}) و سازگاری آسیب‌شناختی \عدد{۹۶٫۵} درصد گزارش شد. تفکیک این دو رقم ضروری است: \emph{۵۰ درصد در انسان، ۷۰ درصد در مدل حیوانی}. این یافته، شاهد دومی است بر همان نکته‌ای که در~\cite{Cho21} دیده شد: \lr{SSIM} بالای \عدد{۰٫۹۸} هیچ تضمینی برای حفظ اطلاعات تشخیصی نمی‌دهد.

آزادبرزین و همکاران~\cite{Aza23} و جو و همکاران~\cite{Jo26} در دو مرور نظام‌مند و گستره‌ای مجزا (به‌ترتیب با \عدد{۱۸} و \عدد{۵۶} مطالعه)، این حوزه را دسته‌بندی کرده‌اند و هر دو بر عدم‌بلوغ گزارش‌دهی کمّی در بخش قابل‌توجهی از این ادبیات تأکید کرده‌اند.

\subsubsection*{شاهد تکمیلی از یک وجهیت مجاور}

در وجهیت‌ها و اندام‌های مجاور نیز همین وظیفه (تولید تصویر با ماده‌ی حاجب از تصویر بدون آن) دنبال شده است، و یک مطالعه از آن‌ها برای منطق ارزیابی این پایان‌نامه اهمیت ویژه دارد. 
فوکودا و همکاران~\cite{Fuk24} با
\lr{CycleGAN} و داده‌ی \عدد{۲۵} بیمار، تباین را در ناحیه‌ی غدد لنفاوی
ژوگولار داخلی هم‌نهادی کردند. \lr{PSNR} تصویر هم‌نهاده (میانه \عدد{۱۷٫۸۰}) از تصویر
بدون‌تزریق خام (\عدد{۱۹٫۷۴}) \emph{بدتر} بود، اما در وظیفه‌ی بالینی تمایز غدد
لنفاوی از عروق، سطح زیر منحنی تصویر هم‌نهاده \عدد{۱٫۰۰} شد، دقیقاً برابر تصویر
با ماده‌ی حاجب واقعی، در حالی که برای تصویر بدون‌تزریق \عدد{۰٫۷۲} بود. مدلی که
بر اساس \lr{PSNR} شکست خورده بود، از نظر بالینی کاملاً موفق بود. این، در کنار
جهش هشت‌برابری نسبت تباین‌به‌نوفه در برابر بهبود \عدد{۹} درصدی \lr{PSNR}
در~\cite{Cho21}، دومین شاهد مستقل بر همان نکته‌ای است که دستورالعمل ارزیابی فصل~\ref{chap:method} بر آن بنا شده است.

\subsubsection*{معماری‌های عمومی به‌عنوان مبنا}

جدا از این مطالعات، گروه دیگری از ادبیات را معماری‌های \emph{عمومی} ترجمه‌ی
تصویر تشکیل می‌دهند: معماری‌هایی که برای هیچ وظیفه‌ی پزشکی خاصی طراحی نشده‌اند
اما به دلیل در دسترس بودن پیاده‌سازی، مبنای متعارف این ادبیات شده‌اند و در
فصل~\ref{chap:method} نیز به‌عنوان مدل مبنا بازآموزی می‌شوند.
\lr{pix2pixHD}~\cite{wang2018pix2pixhd}، \lr{pix2pix} را با یک مولد چندمقیاسی،
متمایزکننده‌های چندمقیاسی و یک \textbf{زیان تطبیق ویژگی} به تفکیک‌پذیری بالا
می‌برد؛ همان زیان تطبیق ویژگی که بعدها در بسیاری از مولدهای پزشکی تکرار شد.
\lr{ResViT}~\cite{dalmaz2022resvit} با بلوک‌های ترادیس‌ساز باقیمانده‌ی تجمیع‌شده،
زمینه‌ی سراسری و جزئیات موضعی را هم‌زمان نگه می‌دارد و از نظر \emph{هدف}
نزدیک‌ترین مبنای منتشرشده به مسئله‌ی حاضر است، هرچند روی \lr{MRI} معرفی شده.
\lr{Ea-GAN}~\cite{yu2019eagans} نقشه‌ی لبه را صریحاً وارد یادگیری می‌کند، در دو
گونه‌ی \lr{gEa-GAN} و \lr{dEa-GAN}؛ یعنی حفظ بافت را از طریق \emph{تابع زیان}
دنبال می‌کند و نه از طریق معیار ارزیابی. سرانجام
\lr{SwinUNETR}~\cite{hatamizadeh2022swinunetr} و
\lr{TransUNet}~\cite{chen2021transunet}، که در اصل برای بخش‌بندی معرفی شده‌اند،
در ادبیات به‌طور گسترده در جای ستون فقرات مولد به کار رفته‌اند.

هیچ‌یک از این معماری‌ها برای سی‌تی چندپویشی معرفی نشده‌اند و هیچ‌کدام سازوکاری
برای شرطی‌سازی بر پویش ندارند؛ پس هر یک باید به‌ازای هر پویش جداگانه آموزش ببیند
و روی محور کنترل‌پذیری اساساً قابل ارزیابی نیست.

\subsubsection*{سرگروه‌های معماری}

پنج خانواده‌ی اصلی بالا، به‌همراه معماری‌های عمومیِ مبنا، از نظر \emph{انگیزه} متفاوت‌اند، اما از نظر \emph{معماری} به شمار کمتری الگو فرو می‌کاهند و تقریباً همه‌ی مطالعات این بخش نمونه‌ای از سه \emph{راهبرد نظارتی} زیرند. جدول~\ref{tab:synthesis-families} هسته‌ی معماری، توانایی اصلی و کاستی اصلی هر خانواده را کنار هم می‌گذارد.

\begin{table}[H]
\centering
\footnotesize
\caption{پنج خانواده‌ی روش‌شناختی هم‌نهادی \lr{NC}-به-\lr{CE}.}
\label{tab:synthesis-families}
\begin{tabular}{|p{2.6cm}|p{3.9cm}|p{3.5cm}|p{3.7cm}|}
\hline
\textbf{خانواده} & \textbf{هسته‌ی معماری} & \textbf{توانایی اصلی} & \textbf{کاستی اصلی} \\
\hline
یکم: \lr{GAN} جفت‌شده & نگاشت شرطی تصویردانه‌به‌تصویردانه با نظارت مستقیم و زیان تخاصمی + بازسازی & بالاترین وفاداری بصری گزارش‌شده؛ قابل هدایت به نواحی مرتبط با ماده‌ی حاجب با توجه یا نقاب & وابسته به جفت هم‌ترازِ هم‌بیمار؛ ناهم‌ترازی باقی‌مانده مستقیماً به خطای آموزش تبدیل می‌شود \\
\hline
دوم: سازگاری چرخه‌ای / بدون‌جفت & دو مولد در دو جهت، با قید بازسازی رفت‌وبرگشت و بدون نیاز به جفت هم‌تراز & آموزش با داده‌ی بدون‌جفت ممکن می‌شود & بدون نظارت تصویردانه‌ای؛ خطر توهم، یعنی تولید ساختار ناموجود \\
\hline
سوم: مدل‌های انتشار & یادگیری فرایند تدریجی حذف نوفه به‌جای آموزش تخاصمی & آموزش پایدارتر و تنوع نمونه‌ی بهتر & استنتاج کند؛ همچنان به هم‌ترازی پیشین وابسته است \\
\hline
چهارم: هم‌نهادی هدایت‌شده با هم‌ترازسازی & هم‌ترازسازی و هم‌نهادی به‌صورت مشترک یا متوالی، برای تفکیک آناتومی مشترک از اثر ماده‌ی حاجب & تنها خانواده‌ای که معیار پایین‌دستی (دایس بخش‌بندی) را مستقیماً گزارش می‌کند & کیفیت هم‌ترازی هرگز به‌صورت کمّی گزارش نشده؛ سهم مدل از سهم داده قابل تفکیک نیست \\
\hline
پنجم: کاربردپذیری بالینی و کاهش دوز & سنجش مستقیم اثر هم‌نهادی بر تصمیم بالینی (حساسیت تشخیص، سازگاری آسیب‌شناختی) به‌جای فقط معیار تصویری & تنها خانواده‌ای که رابطه‌ی معیار تصویردانه‌ای با سودمندی بالینی را مستقیماً می‌آزماید & معیارهای تصویردانه‌ای بالا ضامن حفظ اطلاعات تشخیصی نیستند؛ ارزیابی بالینی پرهزینه و کم‌مقیاس می‌ماند \\
\hline
\end{tabular}
\end{table}

شکل~\ref{fig:syn-archetypes} همین سه راهبرد نظارتی را نشان می‌دهد: \lr{GAN} شرطی جفت‌شده (\lr{pix2pix} و \lr{pix2pixHD}~\cite{wang2018pix2pixhd}) که خانواده‌های یکم، چهارم و پنجم بر آن بنا شده‌اند؛ سازگاری چرخه‌ای (\lr{CycleGAN} و \lr{CyTran}~\cite{Ris21}) که خانواده‌ی دوم است؛ و مدل انتشار شرطی (\lr{CFPS-Diff}~\cite{Zhe25}) که خانواده‌ی سوم است. شکل~\ref{fig:syn-backbone} در کنار آن‌ها ستون فقرات ترکیبی پیچشی--ترادیس‌سازی (\lr{ResViT}~\cite{dalmaz2022resvit}، \lr{SwinUNETR}~\cite{hatamizadeh2022swinunetr}، \lr{TransUNet}~\cite{chen2021transunet}) را نشان می‌دهد که خود راهبرد نظارتی مستقلی نیست بلکه در جای مولدِ هر یک از آن سه می‌نشیند.

تفکیک این الگوها دو نکته را روشن می‌کند. نخست، آن‌چه مطالعات این حوزه را واقعاً از هم متمایز می‌کند اغلب معماری نیست بلکه \emph{راهبرد نظارتی} و \emph{کیفیت داده‌ای} است که در اختیار داشته‌اند؛ افزودن پیمانه توجه یا بلوک ترادیس‌ساز به مولد، در چارچوب همان راهبرد نخست باقی می‌ماند. دوم، وابستگی به هم‌ترازی هندسی، خاصیت راهبرد نخست است و نه خاصیت یک مقاله‌ی خاص: هر روشی که زیان بازسازی تصویردانه‌به‌تصویردانه داشته باشد، ناهم‌ترازی باقی‌مانده‌ی جفت آموزشی را به‌عنوان بخشی از تفاوت \lr{NC}--\lr{CE} می‌آموزد.

% ---------------------------------------------------------------
%  شکل: سه راهبرد نظارتی در سنتز NC-به-CE
%  چیدمان با مختصات صریح (نه positioning نسبی) تا متن‌ها روی هم نیفتند.
%  نیازمند: tikz + arrows.meta
%  استفاده:  % ---------------------------------------------------------------
%  شکل: سه راهبرد نظارتی در سنتز NC-به-CE
%  چیدمان با مختصات صریح (نه positioning نسبی) تا متن‌ها روی هم نیفتند.
%  نیازمند: tikz + arrows.meta
%  استفاده:  % ---------------------------------------------------------------
%  شکل: سه راهبرد نظارتی در سنتز NC-به-CE
%  چیدمان با مختصات صریح (نه positioning نسبی) تا متن‌ها روی هم نیفتند.
%  نیازمند: tikz + arrows.meta
%  استفاده:  \input{figs/arch_synthesis}
% ---------------------------------------------------------------
\begingroup
\providecommand{\ssubtext}[1]{{\color{xtire}\tiny #1}}

\tikzset{
  sb/.style   ={draw=xnokhodi, line width=.5pt, fill=white, rounded corners=1.5pt,
                text width=2.55cm, minimum height=1.05cm, align=center,
                inner xsep=1.2mm, inner ysep=1.2mm, font=\scriptsize},
  snet/.style ={sb, fill=xabikam, draw=xabi},
  sloss/.style={sb, fill=xkamrang},
  sio/.style  ={sb, draw=xkhakestari},
  sflow/.style={->, draw=xtire, line width=.6pt,
                >={Stealth[length=1.7mm,width=1.4mm]}},
  sdash/.style={->, draw=xghermez, line width=.5pt, dash pattern=on 1.8pt off 1.4pt,
                >={Stealth[length=1.7mm,width=1.4mm]}},
  slab/.style ={font=\tiny, color=xkhakestari, inner sep=1.2pt, fill=white},
}

\newcommand{\Shead}[2]{%
  {\small\textbf{#1}}\par\nobreak\vspace{0.4mm}%
  {\scriptsize\color{xtire}#2}\par\nobreak\vspace{2.2mm}}

\begin{figure}[H]
\centering

% ═════════════════ (الف) GAN جفت‌شده ═════════════════
\Shead{(الف) \lr{GAN} شرطی جفت‌شده \,---\, \lr{pix2pix} / \lr{pix2pixHD}}
      {نظارت پیکسل‌به‌پیکسل؛ نیازمند جفت هم‌بیمار و هم‌تراز}
\begin{tikzpicture}
  \node[sio]   (x)  at (     0, 0)     {$x_{\mathrm{NC}}$\\[1pt]\ssubtext{تصویر بدون‌کنتراست}};
  \node[snet]  (g)  at (  -3.5, 0)     {مولد $G$};
  \node[sio]   (yh) at (    -7, 0)     {$\hat{y}_{\mathrm{CE}}$\\[1pt]\ssubtext{خروجی سنتزشده}};
  \node[sloss] (l1) at (    -7,-1.75)  {زیان بازسازی \lr{L1}};
  \node[sio]   (y)  at (    -7,-3.5)   {$y_{\mathrm{CE}}$ واقعی};
  \node[snet]  (d)  at ( -10.6,-1.75)  {متمایزکننده $D$\\[1pt]\ssubtext{زیان تخاصمی}};

  \draw[sflow] (x.west)  -- (g.east);
  \draw[sflow] (g.west)  -- (yh.east);
  \draw[sflow] (yh.south) -- (l1.north);
  \draw[sflow] (y.north)  -- (l1.south);
  \draw[sflow] (yh.west) -| (d.north);
  \draw[sflow] (y.west)  -| (d.south);
  \draw[sdash] (l1.east) -| node[slab, pos=0.35] {پس‌انتشار} (g.south);
\end{tikzpicture}

\par\vspace{7mm}

% ═════════════════ (ب) سازگاری چرخه‌ای ═════════════════
\Shead{(ب) سازگاری چرخه‌ای \,---\, \lr{CycleGAN} / \lr{CyTran}}
      {آموزش با داده‌ی بدون‌جفت؛ بدون نظارت پیکسل‌به‌پیکسل}
\begin{tikzpicture}
  \node[sio]   (x2)  at (     0, 0)    {$x_{\mathrm{NC}}$};
  \node[snet]  (g2)  at (  -3.5, 0)    {$G:\ \mathrm{NC}\!\to\!\mathrm{CE}$};
  \node[sio]   (yh2) at (    -7, 0)    {$\hat{y}_{\mathrm{CE}}$};
  \node[snet]  (d2)  at ( -10.5, 0)    {$D_{\mathrm{CE}}$\\[1pt]\ssubtext{روی مجموعه‌ی \lr{CE}، بدون تناظر یک‌به‌یک}};
  \node[snet]  (f2)  at (    -7,-2.0)  {$F:\ \mathrm{CE}\!\to\!\mathrm{NC}$};
  \node[sio]   (xh2) at (  -3.5,-2.0)  {$\hat{x}_{\mathrm{NC}}$};
  \node[sloss] (c2)  at (     0,-2.0)  {زیان سازگاری چرخه‌ای};

  \draw[sflow] (x2.west)   -- (g2.east);
  \draw[sflow] (g2.west)   -- (yh2.east);
  \draw[sflow] (yh2.west)  -- (d2.east);
  \draw[sflow] (yh2.south) -- (f2.north);
  \draw[sflow] (f2.east)   -- (xh2.west);
  \draw[sflow] (xh2.east)  -- (c2.west);
  \draw[sflow] (x2.south)  -- (c2.north);
\end{tikzpicture}

\par\vspace{7mm}

% ═════════════════ (پ) مدل انتشار شرطی ═════════════════
\Shead{(پ) مدل انتشار شرطی \,---\, \lr{CFPS-Diff}}
      {آموزش پایدارتر از \lr{GAN}؛ استنتاج به‌مراتب کندتر}
\begin{tikzpicture}
  \node[sio]   (y0)  at (     0, 0)    {$y_0 = y_{\mathrm{CE}}$};
  \node[sio]   (yt)  at (  -3.5, 0)    {$y_t$};
  \node[sio]   (yT)  at (    -7, 0)    {$y_T \sim \mathcal{N}(0,I)$};
  \node[sio]   (cnd) at (     0,-2.0)  {$x_{\mathrm{NC}}$\\[1pt]\ssubtext{شرط}};
  \node[snet]  (eps) at (  -3.5,-2.0)  {شبکه‌ی حذف نویز\\[1pt]$\varepsilon_\theta(y_t, t \mid x_{\mathrm{NC}})$};
  \node[sloss] (l3)  at (    -7,-2.0)  {زیان پیش‌بینی نویز};

  \draw[sflow] (y0.west)  -- node[slab, above] {افزودن نویز} (yt.east);
  \draw[sflow] (yt.west)  -- (yT.east);
  \draw[sflow] (yt.south) -- (eps.north);
  \draw[sflow] (cnd.west) -- (eps.east);
  \draw[sflow] (eps.west) -- (l3.east);
  \draw[sdash] (yT.north) -- ++(0,0.7) -- node[slab, above] {استنتاج: $T$ گام حذف نویز معکوس} ++(7,0) -- (y0.north);
\end{tikzpicture}

\caption{سه راهبرد نظارتی در ادبیات سنتز \lr{NC}-به-\lr{CE}. محور تمایز میان (الف) و (ب) وجود یا نبود جفت هم‌بیمارِ هم‌تراز است: (الف) نظارت پیکسل‌به‌پیکسل دارد و به همین دلیل به هم‌ترازی هندسی دقیق وابسته است، و (ب) این نیاز را با قید بازگشت‌پذیری چرخه حذف می‌کند اما تضمینی برای حفظ آناتومی نمی‌دهد. (پ) نظارت را از آموزش تخاصمی به پیش‌بینی نویز منتقل می‌کند و آموزش پایدارتری دارد، اما همچنان به جفت هم‌تراز نیاز دارد. هیچ‌یک از این سه الگو سازوکاری برای \emph{شرطی‌سازی بر سطح تقویت} ندارد.}
\label{fig:syn-archetypes}
\end{figure}
\endgroup

% ---------------------------------------------------------------
\begingroup
\providecommand{\ssubtext}[1]{{\color{xtire}\tiny #1}}

\tikzset{
  sb/.style   ={draw=xnokhodi, line width=.5pt, fill=white, rounded corners=1.5pt,
                text width=2.55cm, minimum height=1.05cm, align=center,
                inner xsep=1.2mm, inner ysep=1.2mm, font=\scriptsize},
  snet/.style ={sb, fill=xabikam, draw=xabi},
  sloss/.style={sb, fill=xkamrang},
  sio/.style  ={sb, draw=xkhakestari},
  sflow/.style={->, draw=xtire, line width=.6pt,
                >={Stealth[length=1.7mm,width=1.4mm]}},
  sdash/.style={->, draw=xghermez, line width=.5pt, dash pattern=on 1.8pt off 1.4pt,
                >={Stealth[length=1.7mm,width=1.4mm]}},
  slab/.style ={font=\tiny, color=xkhakestari, inner sep=1.2pt, fill=white},
}

\newcommand{\Shead}[2]{%
  {\small\textbf{#1}}\par\nobreak\vspace{0.4mm}%
  {\scriptsize\color{xtire}#2}\par\nobreak\vspace{2.2mm}}

\begin{figure}[H]
\centering

% ═════════════════ (الف) GAN جفت‌شده ═════════════════
\Shead{(الف) \lr{GAN} شرطی جفت‌شده \,---\, \lr{pix2pix} / \lr{pix2pixHD}}
      {نظارت پیکسل‌به‌پیکسل؛ نیازمند جفت هم‌بیمار و هم‌تراز}
\begin{tikzpicture}
  \node[sio]   (x)  at (     0, 0)     {$x_{\mathrm{NC}}$\\[1pt]\ssubtext{تصویر بدون‌کنتراست}};
  \node[snet]  (g)  at (  -3.5, 0)     {مولد $G$};
  \node[sio]   (yh) at (    -7, 0)     {$\hat{y}_{\mathrm{CE}}$\\[1pt]\ssubtext{خروجی سنتزشده}};
  \node[sloss] (l1) at (    -7,-1.75)  {زیان بازسازی \lr{L1}};
  \node[sio]   (y)  at (    -7,-3.5)   {$y_{\mathrm{CE}}$ واقعی};
  \node[snet]  (d)  at ( -10.6,-1.75)  {متمایزکننده $D$\\[1pt]\ssubtext{زیان تخاصمی}};

  \draw[sflow] (x.west)  -- (g.east);
  \draw[sflow] (g.west)  -- (yh.east);
  \draw[sflow] (yh.south) -- (l1.north);
  \draw[sflow] (y.north)  -- (l1.south);
  \draw[sflow] (yh.west) -| (d.north);
  \draw[sflow] (y.west)  -| (d.south);
  \draw[sdash] (l1.east) -| node[slab, pos=0.35] {پس‌انتشار} (g.south);
\end{tikzpicture}

\par\vspace{7mm}

% ═════════════════ (ب) سازگاری چرخه‌ای ═════════════════
\Shead{(ب) سازگاری چرخه‌ای \,---\, \lr{CycleGAN} / \lr{CyTran}}
      {آموزش با داده‌ی بدون‌جفت؛ بدون نظارت پیکسل‌به‌پیکسل}
\begin{tikzpicture}
  \node[sio]   (x2)  at (     0, 0)    {$x_{\mathrm{NC}}$};
  \node[snet]  (g2)  at (  -3.5, 0)    {$G:\ \mathrm{NC}\!\to\!\mathrm{CE}$};
  \node[sio]   (yh2) at (    -7, 0)    {$\hat{y}_{\mathrm{CE}}$};
  \node[snet]  (d2)  at ( -10.5, 0)    {$D_{\mathrm{CE}}$\\[1pt]\ssubtext{روی مجموعه‌ی \lr{CE}، بدون تناظر یک‌به‌یک}};
  \node[snet]  (f2)  at (    -7,-2.0)  {$F:\ \mathrm{CE}\!\to\!\mathrm{NC}$};
  \node[sio]   (xh2) at (  -3.5,-2.0)  {$\hat{x}_{\mathrm{NC}}$};
  \node[sloss] (c2)  at (     0,-2.0)  {زیان سازگاری چرخه‌ای};

  \draw[sflow] (x2.west)   -- (g2.east);
  \draw[sflow] (g2.west)   -- (yh2.east);
  \draw[sflow] (yh2.west)  -- (d2.east);
  \draw[sflow] (yh2.south) -- (f2.north);
  \draw[sflow] (f2.east)   -- (xh2.west);
  \draw[sflow] (xh2.east)  -- (c2.west);
  \draw[sflow] (x2.south)  -- (c2.north);
\end{tikzpicture}

\par\vspace{7mm}

% ═════════════════ (پ) مدل انتشار شرطی ═════════════════
\Shead{(پ) مدل انتشار شرطی \,---\, \lr{CFPS-Diff}}
      {آموزش پایدارتر از \lr{GAN}؛ استنتاج به‌مراتب کندتر}
\begin{tikzpicture}
  \node[sio]   (y0)  at (     0, 0)    {$y_0 = y_{\mathrm{CE}}$};
  \node[sio]   (yt)  at (  -3.5, 0)    {$y_t$};
  \node[sio]   (yT)  at (    -7, 0)    {$y_T \sim \mathcal{N}(0,I)$};
  \node[sio]   (cnd) at (     0,-2.0)  {$x_{\mathrm{NC}}$\\[1pt]\ssubtext{شرط}};
  \node[snet]  (eps) at (  -3.5,-2.0)  {شبکه‌ی حذف نویز\\[1pt]$\varepsilon_\theta(y_t, t \mid x_{\mathrm{NC}})$};
  \node[sloss] (l3)  at (    -7,-2.0)  {زیان پیش‌بینی نویز};

  \draw[sflow] (y0.west)  -- node[slab, above] {افزودن نویز} (yt.east);
  \draw[sflow] (yt.west)  -- (yT.east);
  \draw[sflow] (yt.south) -- (eps.north);
  \draw[sflow] (cnd.west) -- (eps.east);
  \draw[sflow] (eps.west) -- (l3.east);
  \draw[sdash] (yT.north) -- ++(0,0.7) -- node[slab, above] {استنتاج: $T$ گام حذف نویز معکوس} ++(7,0) -- (y0.north);
\end{tikzpicture}

\caption{سه راهبرد نظارتی در ادبیات سنتز \lr{NC}-به-\lr{CE}. محور تمایز میان (الف) و (ب) وجود یا نبود جفت هم‌بیمارِ هم‌تراز است: (الف) نظارت پیکسل‌به‌پیکسل دارد و به همین دلیل به هم‌ترازی هندسی دقیق وابسته است، و (ب) این نیاز را با قید بازگشت‌پذیری چرخه حذف می‌کند اما تضمینی برای حفظ آناتومی نمی‌دهد. (پ) نظارت را از آموزش تخاصمی به پیش‌بینی نویز منتقل می‌کند و آموزش پایدارتری دارد، اما همچنان به جفت هم‌تراز نیاز دارد. هیچ‌یک از این سه الگو سازوکاری برای \emph{شرطی‌سازی بر سطح تقویت} ندارد.}
\label{fig:syn-archetypes}
\end{figure}
\endgroup

% ---------------------------------------------------------------
\begingroup
\providecommand{\ssubtext}[1]{{\color{xtire}\tiny #1}}

\tikzset{
  sb/.style   ={draw=xnokhodi, line width=.5pt, fill=white, rounded corners=1.5pt,
                text width=2.55cm, minimum height=1.05cm, align=center,
                inner xsep=1.2mm, inner ysep=1.2mm, font=\scriptsize},
  snet/.style ={sb, fill=xabikam, draw=xabi},
  sloss/.style={sb, fill=xkamrang},
  sio/.style  ={sb, draw=xkhakestari},
  sflow/.style={->, draw=xtire, line width=.6pt,
                >={Stealth[length=1.7mm,width=1.4mm]}},
  sdash/.style={->, draw=xghermez, line width=.5pt, dash pattern=on 1.8pt off 1.4pt,
                >={Stealth[length=1.7mm,width=1.4mm]}},
  slab/.style ={font=\tiny, color=xkhakestari, inner sep=1.2pt, fill=white},
}

\newcommand{\Shead}[2]{%
  {\small\textbf{#1}}\par\nobreak\vspace{0.4mm}%
  {\scriptsize\color{xtire}#2}\par\nobreak\vspace{2.2mm}}

\begin{figure}[H]
\centering

% ═════════════════ (الف) GAN جفت‌شده ═════════════════
\Shead{(الف) \lr{GAN} شرطی جفت‌شده \,---\, \lr{pix2pix} / \lr{pix2pixHD}}
      {نظارت پیکسل‌به‌پیکسل؛ نیازمند جفت هم‌بیمار و هم‌تراز}
\begin{tikzpicture}
  \node[sio]   (x)  at (     0, 0)     {$x_{\mathrm{NC}}$\\[1pt]\ssubtext{تصویر بدون‌کنتراست}};
  \node[snet]  (g)  at (  -3.5, 0)     {مولد $G$};
  \node[sio]   (yh) at (    -7, 0)     {$\hat{y}_{\mathrm{CE}}$\\[1pt]\ssubtext{خروجی سنتزشده}};
  \node[sloss] (l1) at (    -7,-1.75)  {زیان بازسازی \lr{L1}};
  \node[sio]   (y)  at (    -7,-3.5)   {$y_{\mathrm{CE}}$ واقعی};
  \node[snet]  (d)  at ( -10.6,-1.75)  {متمایزکننده $D$\\[1pt]\ssubtext{زیان تخاصمی}};

  \draw[sflow] (x.west)  -- (g.east);
  \draw[sflow] (g.west)  -- (yh.east);
  \draw[sflow] (yh.south) -- (l1.north);
  \draw[sflow] (y.north)  -- (l1.south);
  \draw[sflow] (yh.west) -| (d.north);
  \draw[sflow] (y.west)  -| (d.south);
  \draw[sdash] (l1.east) -| node[slab, pos=0.35] {پس‌انتشار} (g.south);
\end{tikzpicture}

\par\vspace{7mm}

% ═════════════════ (ب) سازگاری چرخه‌ای ═════════════════
\Shead{(ب) سازگاری چرخه‌ای \,---\, \lr{CycleGAN} / \lr{CyTran}}
      {آموزش با داده‌ی بدون‌جفت؛ بدون نظارت پیکسل‌به‌پیکسل}
\begin{tikzpicture}
  \node[sio]   (x2)  at (     0, 0)    {$x_{\mathrm{NC}}$};
  \node[snet]  (g2)  at (  -3.5, 0)    {$G:\ \mathrm{NC}\!\to\!\mathrm{CE}$};
  \node[sio]   (yh2) at (    -7, 0)    {$\hat{y}_{\mathrm{CE}}$};
  \node[snet]  (d2)  at ( -10.5, 0)    {$D_{\mathrm{CE}}$\\[1pt]\ssubtext{روی مجموعه‌ی \lr{CE}، بدون تناظر یک‌به‌یک}};
  \node[snet]  (f2)  at (    -7,-2.0)  {$F:\ \mathrm{CE}\!\to\!\mathrm{NC}$};
  \node[sio]   (xh2) at (  -3.5,-2.0)  {$\hat{x}_{\mathrm{NC}}$};
  \node[sloss] (c2)  at (     0,-2.0)  {زیان سازگاری چرخه‌ای};

  \draw[sflow] (x2.west)   -- (g2.east);
  \draw[sflow] (g2.west)   -- (yh2.east);
  \draw[sflow] (yh2.west)  -- (d2.east);
  \draw[sflow] (yh2.south) -- (f2.north);
  \draw[sflow] (f2.east)   -- (xh2.west);
  \draw[sflow] (xh2.east)  -- (c2.west);
  \draw[sflow] (x2.south)  -- (c2.north);
\end{tikzpicture}

\par\vspace{7mm}

% ═════════════════ (پ) مدل انتشار شرطی ═════════════════
\Shead{(پ) مدل انتشار شرطی \,---\, \lr{CFPS-Diff}}
      {آموزش پایدارتر از \lr{GAN}؛ استنتاج به‌مراتب کندتر}
\begin{tikzpicture}
  \node[sio]   (y0)  at (     0, 0)    {$y_0 = y_{\mathrm{CE}}$};
  \node[sio]   (yt)  at (  -3.5, 0)    {$y_t$};
  \node[sio]   (yT)  at (    -7, 0)    {$y_T \sim \mathcal{N}(0,I)$};
  \node[sio]   (cnd) at (     0,-2.0)  {$x_{\mathrm{NC}}$\\[1pt]\ssubtext{شرط}};
  \node[snet]  (eps) at (  -3.5,-2.0)  {شبکه‌ی حذف نویز\\[1pt]$\varepsilon_\theta(y_t, t \mid x_{\mathrm{NC}})$};
  \node[sloss] (l3)  at (    -7,-2.0)  {زیان پیش‌بینی نویز};

  \draw[sflow] (y0.west)  -- node[slab, above] {افزودن نویز} (yt.east);
  \draw[sflow] (yt.west)  -- (yT.east);
  \draw[sflow] (yt.south) -- (eps.north);
  \draw[sflow] (cnd.west) -- (eps.east);
  \draw[sflow] (eps.west) -- (l3.east);
  \draw[sdash] (yT.north) -- ++(0,0.7) -- node[slab, above] {استنتاج: $T$ گام حذف نویز معکوس} ++(7,0) -- (y0.north);
\end{tikzpicture}

\caption{سه راهبرد نظارتی در ادبیات سنتز \lr{NC}-به-\lr{CE}. محور تمایز میان (الف) و (ب) وجود یا نبود جفت هم‌بیمارِ هم‌تراز است: (الف) نظارت پیکسل‌به‌پیکسل دارد و به همین دلیل به هم‌ترازی هندسی دقیق وابسته است، و (ب) این نیاز را با قید بازگشت‌پذیری چرخه حذف می‌کند اما تضمینی برای حفظ آناتومی نمی‌دهد. (پ) نظارت را از آموزش تخاصمی به پیش‌بینی نویز منتقل می‌کند و آموزش پایدارتری دارد، اما همچنان به جفت هم‌تراز نیاز دارد. هیچ‌یک از این سه الگو سازوکاری برای \emph{شرطی‌سازی بر سطح تقویت} ندارد.}
\label{fig:syn-archetypes}
\end{figure}
\endgroup


% ---------------------------------------------------------------
%  شکل: ستون فقرات ترکیبی پیچشی--ترانسفورمری
%  چیدمان با مختصات صریح (نه positioning نسبی) تا متن‌ها روی هم نیفتند.
%  نیازمند: tikz + arrows.meta
%  استفاده:  % ---------------------------------------------------------------
%  شکل: ستون فقرات ترکیبی پیچشی--ترانسفورمری
%  چیدمان با مختصات صریح (نه positioning نسبی) تا متن‌ها روی هم نیفتند.
%  نیازمند: tikz + arrows.meta
%  استفاده:  % ---------------------------------------------------------------
%  شکل: ستون فقرات ترکیبی پیچشی--ترانسفورمری
%  چیدمان با مختصات صریح (نه positioning نسبی) تا متن‌ها روی هم نیفتند.
%  نیازمند: tikz + arrows.meta
%  استفاده:  \input{figs/arch_synthesis_backbone}
% ---------------------------------------------------------------
\begingroup
\providecommand{\ssubtext}[1]{{\color{xtire}\tiny #1}}

\tikzset{
  sb/.style   ={draw=xnokhodi, line width=.5pt, fill=white, rounded corners=1.5pt,
                text width=2.55cm, minimum height=1.05cm, align=center,
                inner xsep=1.2mm, inner ysep=1.2mm, font=\scriptsize},
  snet/.style ={sb, fill=xabikam, draw=xabi},
  sloss/.style={sb, fill=xkamrang},
  sio/.style  ={sb, draw=xkhakestari},
  sflow/.style={->, draw=xtire, line width=.6pt,
                >={Stealth[length=1.7mm,width=1.4mm]}},
  sdash/.style={->, draw=xghermez, line width=.5pt, dash pattern=on 1.8pt off 1.4pt,
                >={Stealth[length=1.7mm,width=1.4mm]}},
  slab/.style ={font=\tiny, color=xkhakestari, inner sep=1.2pt, fill=white},
}

\newcommand{\Shead}[2]{%
  {\small\textbf{#1}}\par\nobreak\vspace{0.4mm}%
  {\scriptsize\color{xtire}#2}\par\nobreak\vspace{2.2mm}}

\begin{figure}[H]
\centering

% ═════════════════ (ت) ستون فقرات ترکیبی ═════════════════
\Shead{ستون فقرات ترکیبی پیچشی--ترانسفورمری \,---\, \lr{ResViT} / \lr{SwinUNETR} / \lr{TransUNet}}
      {راهبرد نظارتی مستقلی نیست؛ در جای مولد هر یک از سه راهبرد شکل~\ref{fig:syn-archetypes} می‌نشیند}
\begin{tikzpicture}
  \node[sio]  (x4) at (     0, 0)  {$x_{\mathrm{NC}}$};
  \node[snet] (e4) at (  -3.5, 0)  {\beginR رمزگذار پیچشی\endR\\[1pt]\ssubtext{\beginR بافت محلی\endR}};
  \node[snet] (t4) at (    -7, 0)  {\beginR بلوک‌های ترانسفورمر\endR\\[1pt]\ssubtext{\beginR وابستگی دوربرد\endR}};
  \node[snet] (d4) at ( -10.5, 0)  {رمزگشا};

  \draw[sflow] (x4.west) -- (e4.east);
  \draw[sflow] (e4.west) -- (t4.east);
  \draw[sflow] (t4.west) -- (d4.east);
  \draw[sflow] (d4.west) -- ++(-1.2,0)
        node[font=\scriptsize, left, inner sep=1.5pt] {$\hat{y}_{\mathrm{CE}}$};
  \draw[sflow] (e4.north) -- ++(0,0.75)
        -- node[slab, above] {\beginR مسیر پرش\endR} ++(-7,0) -- (d4.north);
\end{tikzpicture}

\caption{ستون فقرات ترکیبی پیچشی--ترانسفورمری. این الگو یک \emph{راهبرد نظارتی} مستقل نیست بلکه ساختار درونی مولد است و می‌تواند جایگزین مولد در هر یک از سه راهبرد شکل~\ref{fig:syn-archetypes} شود. مسیر پیچشی بافت محلی را و مسیر ترانسفورمری وابستگی‌های دوربرد را مدل می‌کند، و مسیر پرش جزئیات مکانی از دست‌رفته در رمزگذاری را بازمی‌گرداند.}
\label{fig:syn-backbone}
\end{figure}
\endgroup

% ---------------------------------------------------------------
\begingroup
\providecommand{\ssubtext}[1]{{\color{xtire}\tiny #1}}

\tikzset{
  sb/.style   ={draw=xnokhodi, line width=.5pt, fill=white, rounded corners=1.5pt,
                text width=2.55cm, minimum height=1.05cm, align=center,
                inner xsep=1.2mm, inner ysep=1.2mm, font=\scriptsize},
  snet/.style ={sb, fill=xabikam, draw=xabi},
  sloss/.style={sb, fill=xkamrang},
  sio/.style  ={sb, draw=xkhakestari},
  sflow/.style={->, draw=xtire, line width=.6pt,
                >={Stealth[length=1.7mm,width=1.4mm]}},
  sdash/.style={->, draw=xghermez, line width=.5pt, dash pattern=on 1.8pt off 1.4pt,
                >={Stealth[length=1.7mm,width=1.4mm]}},
  slab/.style ={font=\tiny, color=xkhakestari, inner sep=1.2pt, fill=white},
}

\newcommand{\Shead}[2]{%
  {\small\textbf{#1}}\par\nobreak\vspace{0.4mm}%
  {\scriptsize\color{xtire}#2}\par\nobreak\vspace{2.2mm}}

\begin{figure}[H]
\centering

% ═════════════════ (ت) ستون فقرات ترکیبی ═════════════════
\Shead{ستون فقرات ترکیبی پیچشی--ترانسفورمری \,---\, \lr{ResViT} / \lr{SwinUNETR} / \lr{TransUNet}}
      {راهبرد نظارتی مستقلی نیست؛ در جای مولد هر یک از سه راهبرد شکل~\ref{fig:syn-archetypes} می‌نشیند}
\begin{tikzpicture}
  \node[sio]  (x4) at (     0, 0)  {$x_{\mathrm{NC}}$};
  \node[snet] (e4) at (  -3.5, 0)  {\beginR رمزگذار پیچشی\endR\\[1pt]\ssubtext{\beginR بافت محلی\endR}};
  \node[snet] (t4) at (    -7, 0)  {\beginR بلوک‌های ترانسفورمر\endR\\[1pt]\ssubtext{\beginR وابستگی دوربرد\endR}};
  \node[snet] (d4) at ( -10.5, 0)  {رمزگشا};

  \draw[sflow] (x4.west) -- (e4.east);
  \draw[sflow] (e4.west) -- (t4.east);
  \draw[sflow] (t4.west) -- (d4.east);
  \draw[sflow] (d4.west) -- ++(-1.2,0)
        node[font=\scriptsize, left, inner sep=1.5pt] {$\hat{y}_{\mathrm{CE}}$};
  \draw[sflow] (e4.north) -- ++(0,0.75)
        -- node[slab, above] {\beginR مسیر پرش\endR} ++(-7,0) -- (d4.north);
\end{tikzpicture}

\caption{ستون فقرات ترکیبی پیچشی--ترانسفورمری. این الگو یک \emph{راهبرد نظارتی} مستقل نیست بلکه ساختار درونی مولد است و می‌تواند جایگزین مولد در هر یک از سه راهبرد شکل~\ref{fig:syn-archetypes} شود. مسیر پیچشی بافت محلی را و مسیر ترانسفورمری وابستگی‌های دوربرد را مدل می‌کند، و مسیر پرش جزئیات مکانی از دست‌رفته در رمزگذاری را بازمی‌گرداند.}
\label{fig:syn-backbone}
\end{figure}
\endgroup

% ---------------------------------------------------------------
\begingroup
\providecommand{\ssubtext}[1]{{\color{xtire}\tiny #1}}

\tikzset{
  sb/.style   ={draw=xnokhodi, line width=.5pt, fill=white, rounded corners=1.5pt,
                text width=2.55cm, minimum height=1.05cm, align=center,
                inner xsep=1.2mm, inner ysep=1.2mm, font=\scriptsize},
  snet/.style ={sb, fill=xabikam, draw=xabi},
  sloss/.style={sb, fill=xkamrang},
  sio/.style  ={sb, draw=xkhakestari},
  sflow/.style={->, draw=xtire, line width=.6pt,
                >={Stealth[length=1.7mm,width=1.4mm]}},
  sdash/.style={->, draw=xghermez, line width=.5pt, dash pattern=on 1.8pt off 1.4pt,
                >={Stealth[length=1.7mm,width=1.4mm]}},
  slab/.style ={font=\tiny, color=xkhakestari, inner sep=1.2pt, fill=white},
}

\newcommand{\Shead}[2]{%
  {\small\textbf{#1}}\par\nobreak\vspace{0.4mm}%
  {\scriptsize\color{xtire}#2}\par\nobreak\vspace{2.2mm}}

\begin{figure}[H]
\centering

% ═════════════════ (ت) ستون فقرات ترکیبی ═════════════════
\Shead{ستون فقرات ترکیبی پیچشی--ترانسفورمری \,---\, \lr{ResViT} / \lr{SwinUNETR} / \lr{TransUNet}}
      {راهبرد نظارتی مستقلی نیست؛ در جای مولد هر یک از سه راهبرد شکل~\ref{fig:syn-archetypes} می‌نشیند}
\begin{tikzpicture}
  \node[sio]  (x4) at (     0, 0)  {$x_{\mathrm{NC}}$};
  \node[snet] (e4) at (  -3.5, 0)  {\beginR رمزگذار پیچشی\endR\\[1pt]\ssubtext{\beginR بافت محلی\endR}};
  \node[snet] (t4) at (    -7, 0)  {\beginR بلوک‌های ترانسفورمر\endR\\[1pt]\ssubtext{\beginR وابستگی دوربرد\endR}};
  \node[snet] (d4) at ( -10.5, 0)  {رمزگشا};

  \draw[sflow] (x4.west) -- (e4.east);
  \draw[sflow] (e4.west) -- (t4.east);
  \draw[sflow] (t4.west) -- (d4.east);
  \draw[sflow] (d4.west) -- ++(-1.2,0)
        node[font=\scriptsize, left, inner sep=1.5pt] {$\hat{y}_{\mathrm{CE}}$};
  \draw[sflow] (e4.north) -- ++(0,0.75)
        -- node[slab, above] {\beginR مسیر پرش\endR} ++(-7,0) -- (d4.north);
\end{tikzpicture}

\caption{ستون فقرات ترکیبی پیچشی--ترانسفورمری. این الگو یک \emph{راهبرد نظارتی} مستقل نیست بلکه ساختار درونی مولد است و می‌تواند جایگزین مولد در هر یک از سه راهبرد شکل~\ref{fig:syn-archetypes} شود. مسیر پیچشی بافت محلی را و مسیر ترانسفورمری وابستگی‌های دوربرد را مدل می‌کند، و مسیر پرش جزئیات مکانی از دست‌رفته در رمزگذاری را بازمی‌گرداند.}
\label{fig:syn-backbone}
\end{figure}
\endgroup


\subsubsection*{جمع‌بندی}

این مرور سه نکته‌ی مشترک را آشکار می‌کند. نخست، هیچ‌یک از این مطالعات معیاری برای سنجش کمّی کیفیت هم‌ترازی داده‌ی آموزشی گزارش نکرده‌اند. بدون چنین معیاری، تفکیک سهم کیفیت داده از سهم خود مدل در نتیجه‌ی نهایی ممکن نیست: نمی‌توان تشخیص داد مدل کمبود کیفیت هم‌ترازی را جبران کرده یا نتیجه صرفاً بازتاب داده‌ای است که از پیش خوب هم‌تراز شده بود.

دوم، \emph{ناهمگنی معیارها}. حتی مطالعاتی که هر دو معیار \lr{PSNR} و \lr{SSIM} را گزارش می‌کنند، ارقامشان در بازه‌های کاملاً متفاوتی قرار دارد: از \lr{PSNR} برابر \عدد{۱۷٫۲۳} در~\cite{Liu22} تا \عدد{۴۳٫۲۵} در~\cite{Han24} و حتی \عدد{۵۷} در~\cite{Pin23}؛ اختلافی بیش از \عدد{۴۰} دسی‌بل. این اختلاف عمدتاً بازتاب تفاوت در پنجره‌ی شدت روشنایی، دامنه‌ی محاسبه (کل تصویر در برابر درون نقاب بدن یا نواحی کوچک تشریحی) و نوع داده است و نه تفاوت واقعی کیفیت؛ پس مقایسه‌ی مستقیم این اعداد میان سطرها بی‌معناست و این ادبیات به یک دستورالعمل ارزیابی یکدست نیاز دارد. همچنین، از میان مطالعات چهار خانواده‌ی نخست جدول~\ref{tab:synthesis-summary} تنها سه مورد کد تأییدشده دارند (\cite{Cho21} روی قفسه‌سینه، \cite{Ris21} روی ریه و \cite{Pin23} روی سی‌تی مداخله‌ای کلیوی) و هیچ‌یک به هم‌نهادی چندپویشی سی‌تی \emph{شکمی تشخیصی} نمی‌پردازد.

سوم، \emph{ابهام روش شرطی‌سازی بر پویش}. عبارت «هر دو پویش» در ستون «پویش داده‌ی آموزش» تنها نشان می‌دهد داده‌ی آموزش دو پویش را پوشش می‌دهد و چیزی درباره‌ی نحوه‌ی تولید خروجی نمی‌گوید. بیشتر مطالعات نتیجه‌ی هر پویش را جداگانه گزارش می‌کنند، بی‌آن‌که روشن کنند این ارقام از دو مدل مستقل (با معماری یکسان برای هر پویش مقصد) آمده‌اند یا از یک شبکه‌ی بدون نشانگر پویش؛ مقالات اصلی این دو حالت را از هم تفکیک نمی‌کنند. استثنای مستند پینوک و همکاران~\cite{Pin23} است که یک شبکه‌ی واحد را با نشانگر پویش شرطی کردند و نشان دادند نتیجه با آموزش دو مدل جداگانه هم‌رده است، هرچند دامنه‌ی آزمونشان سی‌تی مداخله‌ای کلیه است و نه هم‌نهادی شکمی تشخیصی. مورد محتمل دیگر، که پیش‌تر در همین فصل معرفی شد، مدل انتشار شرطی \lr{CFPS-Diff}~\cite{Zhe25} است که پویایی گذر ماده‌ی حاجب را \emph{میان} پویش‌ها مدل می‌کند؛ اما در مقاله روشن نیست این مدل بر پویش هدف شرطی شده یا برای هر پویش جداگانه اجرا می‌شود، و در این پایان‌نامه آزموده نشده است.

جدول~\ref{tab:synthesis-summary} این ادبیات را بر اساس پنج خانواده مرتب می‌کند و برای هر مطالعه مجموعه‌داده، مقادیر \lr{PSNR} و \lr{SSIM}، سایر معیارهای گزارش‌شده و وضعیت انتشار کد را کنار هم می‌گذارد. جدول تنها مطالعاتی را در بر می‌گیرد که مقدار کمّی قابل‌مقایسه‌ای گزارش کرده‌اند؛ خط تیره در ستون‌های وفاداری تصویری یعنی ارزیابی مطالعه کاملاً پایین‌دستی یا بالینی است.

{\scriptsize
\begin{longtable}{|p{2.05cm}|p{0.45cm}|p{2.35cm}|p{1.15cm}|p{1.05cm}|p{1.05cm}|p{2.3cm}|p{0.45cm}|}
\caption{ادبیات هم‌نهادی \lr{NC}-به-\lr{CE}، دسته‌بندی‌شده بر اساس خانواده‌ی روش. ستون «پویش داده‌ی آموزش» نشان می‌دهد آیا داده‌ی آموزش هر دو پویش را به‌صورت جفت‌شده در بر داشته است یا نه.
\lr{PSNR} برحسب دسی‌بل، همه‌ی مقادیر \lr{SSIM}، \lr{Dice}، دقت، حساسیت و \lr{AUC} در بازه‌ی صفر تا یک، و \lr{MAE} و \lr{HD} برحسب واحد \lr{HU} و میلی‌متر؛ \checkmark یعنی کد تأییدشده در دسترس است.
خط تیره در ستون‌های \lr{PSNR} و \lr{SSIM} یعنی مطالعه معیار وفاداری تصویری را به‌عنوان نتیجه‌ی اصلی گزارش نمی‌کند و ارزیابی‌اش کاملاً پایین‌دستی است.
مقدار \lr{SSIM} سطر \lr{pix2pix} سه‌بعدی در واقع \lr{MS-SSIM} است (متن بالا)؛ \lr{CyTran} به‌جای \lr{PSNR}، \lr{MAE} و \lr{RMSE} گزارش می‌کند و مبنای بدون‌ترجمه در همان جدول \lr{SSIM} برابر \عدد{۰٫۶۵۶} دارد.
در سطر یادگیری چندوظیفه‌ای متحد، ستون‌های \lr{PSNR} و \lr{SSIM} ارقام وضوح بالا را نشان می‌دهند و ارقام وضوح پایین در ستون سایر معیارها آمده است.
\textsuperscript{\dag}\lr{Res-CycleGAN} در خانواده‌ی دوم (سازگاری چرخه‌ای) آمده چون از معماری \lr{CycleGAN} استفاده می‌کند، اما برخلاف تعریف آن خانواده، خودِ مقاله تصریح می‌کند داده‌ی آموزش \emph{جفت‌شده} بوده است.
\textsuperscript{\ddag}\lr{Pin23}: ستون \lr{Dice} مرجع تنها مدل شرطی‌شده را نشان می‌دهد (متن بالا)؛ برای مرجع، دایسِ بخش‌بندی روی تصویر \emph{واقعی} (نه هم‌نهاده) \عدد{۰٫۴۴} (\lr{CME}) و \عدد{۰٫۶۵} (\lr{NGE}) بود، یعنی برای پویش \lr{CME} حتی از دایس تصویر هم‌نهاده کمتر. \lr{PSNR} این مطالعه در مقیاسی به‌مراتب بالاتر از بقیه‌ی جدول است و مستقیماً با آن قابل‌مقایسه نیست.}
\label{tab:synthesis-summary} \\
\hline
\textbf{روش (مرجع)} & \textbf{سال} & \textbf{داده و وظیفه} & \textbf{پویش داده‌ی آموزش} & \textbf{\lr{PSNR}} & \textbf{\lr{SSIM}} & \textbf{سایر معیارها} & \textbf{کد} \\
\hline
\endfirsthead
\multicolumn{8}{r}{\scriptsize \emph{ادامه‌ی جدول~\ref{tab:synthesis-summary} از صفحه‌ی پیش}} \\
\hline
\textbf{روش (مرجع)} & \textbf{سال} & \textbf{داده و وظیفه} & \textbf{پویش داده‌ی آموزش} & \textbf{\lr{PSNR}} & \textbf{\lr{SSIM}} & \textbf{سایر معیارها} & \textbf{کد} \\
\hline
\endhead
\hline
\multicolumn{8}{r}{\scriptsize \emph{ادامه در صفحه‌ی بعد}} \\
\endfoot
\endlastfoot
\multicolumn{8}{|r|}{\textbf{خانواده‌ی یکم: \lr{GAN} جفت‌شده\quad\lr{(Paired GAN)}}} \\
\hline
\lr{Dual-path GAN}~\cite{Liu22} & \عدد{۲۰۲۲} & خصوصی (\عدد{۶۵} بیمار)؛ هم‌نهادی سه پویش شکم و لگن & هر دو پویش & شریانی: \عدد{۱۹٫۰۹} \newline وریدی: \عدد{۱۹٫۷۹} \newline تأخیری: \عدد{۱۷٫۲۳} & شریانی: \عدد{۰٫۶۳} \newline وریدی: \عدد{۰٫۶۵} \newline تأخیری: \عدد{۰٫۶۵} & --- & \xmark \\
\hline
\lr{cGAN} بهبودیافته~\cite{Yan22} & \عدد{۲۰۲۲} & خصوصی (\عدد{۸۶} بیمار)؛ هم‌نهادی شریانی و وریدی از ضایعات کبدی & هر دو پویش & شریانی: \عدد{۳۳٫۱۸} \newline وریدی: \عدد{۳۳٫۰۲} & شریانی: \عدد{۰٫۸۱۹} \newline وریدی: \عدد{۰٫۸۱۹} & \lr{PCC} \عدد{۰٫۹۵۶}؛ \عدد{۰٫۹۶۱} & \xmark \\
\hline
\lr{SAG-GAN}~\cite{Yan23} & \عدد{۲۰۲۳} & همان داده‌ی بالا؛ افزودن توجه فضایی بدون‌نظارت & هر دو پویش & شریانی: \عدد{۳۳٫۱۸} \newline وریدی: \عدد{۳۳٫۰۵} & شریانی: \عدد{۰٫۸۳۰} \newline وریدی: \عدد{۰٫۸۲۳} & \lr{PCC} \عدد{۰٫۹۶۶}؛ \عدد{۰٫۹۶۸} \newline \lr{SSIM} کبد \عدد{۰٫۹۳۲} & \xmark \\
\hline
\lr{SGCDD-GAN}~\cite{Yan24c} & \عدد{۲۰۲۴} & همان داده؛ رمزگشایی دوگانه + هدایت بخش‌بندی & هر دو پویش & شریانی: \عدد{۳۳٫۳} \newline وریدی: \عدد{۳۳٫۱۴} & شریانی: \عدد{۰٫۸۳۷۱} \newline وریدی: \عدد{۰٫۸۲۴۷} & دقت دسته‌بندی ضایعه \عدد{۰٫۸۲۷} \newline دقت رادیولوژیست \عدد{۰٫۹۴۱} \newline \lr{MSE} شریانی \عدد{۳۱٫۰۷}، وریدی \عدد{۳۲٫۱۶} \newline \lr{PCC} شریانی \عدد{۰٫۹۷۷۰}، وریدی \عدد{۰٫۹۷۸۸} & \xmark \\
\hline
\lr{MAN-GAN}~\cite{Zha25} & \عدد{۲۰۲۵} & \عدد{۳۷۴} آموزش / \عدد{۱۴۴} آزمون داخلی / \عدد{۸۳} خارجی؛ سه پویش کبدی & هر دو پویش & شریانی: \عدد{۲۷٫۷۱} \newline وریدی: \عدد{۲۶٫۷۶} \newline تأخیری: \عدد{۲۶٫۱۷} & شریانی: \عدد{۰٫۸۲} \newline وریدی: \عدد{۰٫۸۰} \newline تأخیری: \عدد{۰٫۷۹} & مبنای \lr{NC} خام: \عدد{۲۶٫۰۲}/\عدد{۰٫۷۴} \newline ارزیابی رادیولوژیست & \xmark \\
\hline
\lr{MAFormer}~\cite{Che24} & \عدد{۲۰۲۴} & خصوصی (\عدد{۳۷۴} بیمار)؛ ترجمه‌ی چندپویشی کبدی با تخمین تصویر تفریق & هر دو پویش & شریانی: \عدد{۲۸٫۰۱} \newline وریدی: \عدد{۲۶٫۷۹} \newline تأخیری: \عدد{۲۶٫۱۹} & شریانی: \عدد{۰٫۸۲۲} \newline وریدی: \عدد{۰٫۷۹۷} \newline تأخیری: \عدد{۰٫۷۸۸} & --- & \xmark \\
\hline
\lr{AMNet}~\cite{Zhu25c} & \عدد{۲۰۲۵} & \عدد{۹۲} حجم (کلینیک کلیولند)؛ هم‌نهادی پویش وریدی & هر دو پویش (هم‌ترازشده) & کل‌نگر: \عدد{۲۴٫۸۸} \newline تومور: \عدد{۲۵٫۲۰} & کل‌نگر: \عدد{۰٫۸۵} \newline تومور: \عدد{۰٫۸۵} & \lr{MAE} کل‌نگر \عدد{۶٫۲۸} \newline \lr{MAE} تومور \عدد{۶٫۱۷} & \xmark \\
\hline
\lr{pix2pix} سه‌بعدی~\cite{Cho21} & \عدد{۲۰۲۱} & \عدد{۲۵} توسعه / \عدد{۲۵} آزمون فنی / \عدد{۱۲} آزمون بالینی؛ قفسه‌سینه & هر دو پویش & هم‌نهادی: \عدد{۱۷٫۴۴} \newline مبنا \lr{NC}: \عدد{۱۵٫۹۷} & هم‌نهادی: \عدد{۰٫۸۴} \newline مبنا \lr{NC}: \عدد{۰٫۸۱} & \lr{MAE} \عدد{۴۱٫۷۲}؛ \lr{LPIPS} \عدد{۰٫۱۴} \newline \lr{CNR} \عدد{۶٫۱۵} در برابر \عدد{۰٫۷۴} & \checkmark \\
\hline
\lr{Phase-conditioned UNet}~\cite{Pin23} & \عدد{۲۰۲۳} & سی‌تی مداخله‌ای؛ هدایت کرایوتراپی کلیوی & هر دو پویش، \emph{شرطی‌شده}\textsuperscript{\ddag} & \lr{CME}: \عدد{۵۷٫۰۹} \newline \lr{NGE}: \عدد{۵۷٫۰۴} & \lr{CME}: \عدد{۰٫۹۹۴۳} \newline \lr{NGE}: \عدد{۰٫۹۹۳۶} & \lr{Dice} تومور: \lr{CME} \عدد{۰٫۵۰}، \lr{NGE} \عدد{۰٫۴۴} & \checkmark \\
\hline
\multicolumn{8}{|r|}{\textbf{خانواده‌ی دوم: سازگاری چرخه‌ای / آموزش بدون‌جفت}} \\
\hline
\lr{Res-CycleGAN}~\cite{Xie21} & \عدد{۲۰۲۱} & \عدد{۲۰} بیمار کبد؛ سی‌تی برنامه‌ریزی پرتودرمانی شکم & جفت‌شده\textsuperscript{\dag} & هم‌نهادی: \عدد{۲۶٫۵} \newline مبنا \lr{NC}: \عدد{۲۲٫۷} & هم‌نهادی: \عدد{۰٫۵۵۸} \newline مبنا \lr{NC}: \عدد{۰٫۳۳۱} & \lr{MAE} هم‌نهادی \عدد{۴۶٫۲}؛ مبنا \عدد{۵۹٫۳} & \xmark \\
\hline
\lr{CyTran}~\cite{Ris21} & \عدد{۲۰۲۱} & \lr{Coltea-Lung-CT-100W} (عمومی)؛ \lr{NC}-به-\lr{CE} ریه & بدون‌جفت & گزارش‌نشده & وریدی: \عدد{۰٫۷۴۵} \newline شریانی: \عدد{۰٫۷۵۸} & \lr{MAE} \عدد{۰٫۰۶۱}؛ \عدد{۰٫۰۵۹} \newline \lr{RMSE} \عدد{۰٫۱۴۴}؛ \عدد{۰٫۱۴۷} & \checkmark \\
\hline
\multicolumn{8}{|r|}{\textbf{خانواده‌ی سوم: مدل‌های انتشار}} \\
\hline
\lr{CFPS-Diff}~\cite{Zhe25} & \عدد{۲۰۲۵} & \عدد{۵۴۷} آموزش (یک مرکز) / \عدد{۱۱۶} آزمون (دو مرکز) / \عدد{۲۸} خارجی؛ هم‌نهادی سه پویش & هر دو پویش (هم‌ترازشده) & آزمون: \عدد{۲۹٫۳۹} \newline خارجی: \عدد{۲۹٫۴۲} & آزمون: \عدد{۰٫۸۲۵} \newline خارجی: \عدد{۰٫۸۲۶} & \lr{MAE} \عدد{۳۱٫۲۰}؛ \عدد{۳۰٫۷۵} \newline \lr{FID}، \lr{LPIPS} گزارش‌شده & \xmark \\
\hline
\multicolumn{8}{|r|}{\textbf{خانواده‌ی چهارم: هم‌نهادی هدایت‌شده با هم‌ترازسازی تصویر}} \\
\hline
\lr{CKAP-Net}~\cite{Zho24} & \عدد{۲۰۲۴} & \عدد{۴۰۳} + \عدد{۳۸} بیمار (دو مرکز)؛ هم‌نهادی شریانی دیررس + بخش‌بندی \عدد{۱۳} اندام & هر دو پویش (جفت‌شده) & \عدد{۳۲٫۶۰} \newline (مبنا: \عدد{۲۹٫۲۸}) & \عدد{۰٫۸۵۵} \newline (مبنا: \عدد{۰٫۷۸۰}) & \lr{MAE} \عدد{۲۵٫۹۶} \newline \lr{Dice} پایین‌دستی \عدد{۰٫۸۱۲} \newline \lr{HD} \عدد{۵٫۴۶} & \xmark \\
\hline
یادگیری چندوظیفه‌ای متحد~\cite{Zho23} & \عدد{۲۰۲۳} & شکمی؛ هم‌ترازسازی تغییرشکل‌پذیر و هم‌نهادی به‌صورت مشترک & هر دو پویش & شریانی: \عدد{۲۶٫۲۷} \newline وریدی: \عدد{۲۸٫۲۵} & شریانی: \عدد{۰٫۹۲} \newline وریدی: \عدد{۰٫۹۶} & \lr{MAE} شریانی \عدد{۴۱٫۸۹}، وریدی \عدد{۳۱٫۱۸} \newline وضوح پایین: \lr{PSNR} \عدد{۲۷٫۶۹} و \عدد{۲۷٫۴۴}، \lr{SSIM} \عدد{۰٫۹۵} و \عدد{۰٫۹۶}، \lr{MAE} \عدد{۲۸٫۶۶} و \عدد{۳۱٫۷۱} & \xmark \\
\hline
تفکیک سازگاری آناتومیک~\cite{Li24} & \عدد{۲۰۲۴} & \عدد{۲۷۸} بیمار سرطان کبد + \lr{SynthRAD2023} لگنی؛ جهت \lr{CECT}-به-\lr{CT} & هر دو پویش (هم‌ترازشده) & \عدد{۳۶٫۵۲} & \عدد{۰٫۸۹۲} & \lr{MAE} \عدد{۱۶٫۳۸} & \xmark \\
\hline
چارچوب متحد کلیوی~\cite{Uhm22} & \عدد{۲۰۲۲} & سی‌تی چندپویشی کلیه؛ هم‌ترازسازی \lr{DEEDS} پیش از هم‌نهادی + دسته‌بندی زیرگونه & هر دو پویش (هم‌ترازشده) & --- & --- & \lr{mAUC} زیرگونه \عدد{۰٫۸۳۴} \newline (سه‌پویش \عدد{۰٫۸۱۵}؛ سقف چهارپویش \عدد{۰٫۸۳۵}) & \xmark \\
\hline
\multicolumn{8}{|r|}{\textbf{خانواده‌ی پنجم: کاربردپذیری بالینی و کاهش دوز ماده‌ی حاجب}} \\
\hline
\lr{DL-SynCCT}~\cite{Han24} & \عدد{۲۰۲۴} & \عدد{۲٬۲۰۲} جفت آموزش (\lr{DECT})؛ \عدد{۳۹۸} بیمار خارجی از سه مرکز & هر دو پویش (\lr{VNC}) & \عدد{۴۳٫۲۵} & \عدد{۰٫۹۲} & حساسیت تشخیص ضایعه \عدد{۰٫۷۲۰} به \عدد{۰٫۷۶۴} \newline اطمینان \عدد{۳٫۰} به \عدد{۳٫۶} & \xmark \\
\hline
\lr{cGAN} دومرحله‌ای~\cite{Kim21} & \عدد{۲۰۲۱} & \عدد{۳۵۳} بیمار اورژانس از سه مرکز؛ مطالعه‌ی خوانندگان با شش رادیولوژیست & هر دو پویش & --- & --- & دقت تشخیصی: بهبود در \عدد{۳} از \عدد{۶} خواننده \newline اطمینان: بهبود در \عدد{۵} از \عدد{۶} ($p<0.001$) & \xmark \\
\hline
\lr{cGAN} کاهش دوز~\cite{Hau21} & \عدد{۲۰۲۱} & \عدد{۸۵} بیمار (\lr{DECT}، پویش شریانی)؛ بازسازی دوز کامل از دوز کاهش‌یافته & هر دو پویش (\lr{VNC}) & کاهش \عدد{۵۰}٪: $>$\عدد{۵۱} \newline کاهش \عدد{۸۰}٪: $>$\عدد{۴۸} & کاهش \عدد{۵۰}٪: $>$\عدد{۰٫۹۹} \newline کاهش \عدد{۸۰}٪: $>$\عدد{۰٫۹۸} & \lr{FID} \عدد{۰٫۸}--\عدد{۱٫۷} \newline سازگاری آسیب‌شناختی \عدد{۱۰۰}٪ فقط در کاهش \عدد{۵۰}٪ & \xmark \\
\hline
مطالعه‌ی حیوانی~\cite{Hau22} & \عدد{۲۰۲۲} & \عدد{۲۰} مینی‌پیگ؛ کاهش دوز ماده‌ی حاجب & هر دو پویش & --- & --- & \lr{CNR} عروقی \عدد{۶٫۹} $\to$ \عدد{۲۸٫۷} (واقعی \عدد{۳۴٫۲}) \newline سازگاری آسیب‌شناختی \عدد{۹۶٫۵}٪ \newline کاهش دوز \عدد{۷۰}٪ & \xmark \\
\hline
\end{longtable}
}


\section{مقایسه‌ی دسته‌ها و جای خالی پژوهش}
\label{sec:gap-summary}

سه حوزه‌ی مرورشده سه پاسخ متفاوت به یک مشاهده‌ی واحد‌اند: مدل‌های بخش‌بندی روی پویش بدون‌تزریق کم‌دقت‌ترند. جدول~\ref{tab:lit-overview} این سه دسته را کنار هم می‌گذارد و جایگاه این پایان‌نامه را نسبت به هر یک نشان می‌دهد.

\begin{table}[H]
\centering
\footnotesize
\caption{مقایسه‌ی سه دسته‌ی رویکرد به شکاف بخش‌بندی در پویش بدون‌تزریق و جایگاه این پایان‌نامه.}
\label{tab:lit-overview}
\begin{tabular}{|p{2.0cm}|p{2.7cm}|p{2.7cm}|p{2.9cm}|p{3.0cm}|}
\hline
\textbf{دسته} & \textbf{ایده} & \textbf{توانایی اصلی} & \textbf{کاستی اصلی} & \textbf{جایگاه این پایان‌نامه} \\
\hline
آموزش مستقیم روی پویش بدون‌تزریق~\cite{park2025nect} & مدل بخش‌بندی روی داده‌ی برچسب‌خورده‌ی بدون‌تزریق آموزش می‌بیند & در اندام‌های بزرگ و جمعیت همگن به سطح پویش با ماده‌ی حاجب می‌رسد & به داده‌ی برچسب‌خورده‌ی بومی هر مرکز نیاز دارد؛ اندام‌های دشوار مانند پانکراس پوشش داده نشده‌اند & جایگزین نمی‌شود؛ داده‌ی لازم در دسترس نیست (فصل~\ref{chap:survey}) \\
\hline
هم‌ترازسازی بین‌پویشی & پویش‌های موجود یک بیمار هندسی بر هم منطبق می‌شوند & جفت‌های هم‌ترازِ لازم برای هم‌نهادی را می‌سازد & معیارهای شدت‌محور در برابر تباین شکننده‌اند؛ روش‌های آگاه از آناتومی تنها وزن‌دهی می‌کنند & زنجیره‌ی صلب که دامنه‌ی معیار را به اندام‌های پایدار محدود می‌کند؛ گام ساخت داده \\
\hline
هم‌نهادی \lr{NC}-به-\lr{CE} & پویش با ماده‌ی حاجب مستقیماً از پویش بدون‌تزریق بازسازی می‌شود & مسئله‌ی سرتاسری؛ اثر بالینی و پایین‌دستی گزارش‌شده & وابسته به کیفیت هم‌ترازی؛ ارزیابی ناهمگن؛ تقریباً بدون شرط‌سازی بر پویش و سطح تقویت & \textbf{هسته‌ی پایان‌نامه:} \lr{PLC-Net} با شرط پویش و سطح و ارزیابی سطح‌اندام و پایین‌دستی \\
\hline
\end{tabular}
\end{table}

شکاف بخش‌بندی \emph{وجود دارد} اما یکنواخت نیست: در تنها اندازه‌گیری کنترل‌شده‌ی موجود~\cite{jeon2024fully}، این افت برای اندام‌های بزرگ و پرتباین اندک است و برای اندامی کوچک و کم‌مرز مانند پانکراس به‌روشنی بزرگ‌تر می‌شود. شکاف \emph{تماماً ذاتی نیز نیست}: آموزش مستقیم روی داده‌ی بدون‌تزریق، دست‌کم برای اندام‌های بزرگ توپر، بخش عمده‌ی آن را می‌بندد~\cite{park2025nect}؛ اما این راه به داده‌ای نیاز دارد که فصل~\ref{chap:survey} نشان داد به‌اندازه‌ی کافی وجود ندارد. ادبیات از همین رو دو راهبرد مکمل دارد.

در هم‌ترازسازی، معیارهای کلاسیک در برابر تغییر ناهمگن تباین شکننده‌اند و مسئله‌ی هم‌ترازسازی \emph{بین‌پویشی} شکمی کم بررسی شده است؛ تنها کار موجود~\cite{park2021multiphase} است. نکته‌ی مشترک \emph{همه‌ی} روش‌ها، از آن‌ها که آناتومی را برای وزن‌دهی به کار می‌گیرند تا آن‌ها که معیار را مقاوم‌تر می‌سازند، آن است که معیار تراز را همچنان روی \emph{کل} میدان دید محاسبه می‌کنند.

در هم‌نهادی، از \lr{GAN}های جفت‌شده تا مدل‌های انتشار اخیر، بازسازی مستقیم اطلاعات تباین ممکن است، اما کیفیت آن همچنان به دقت هم‌ترازی پیشین وابسته است. تقریباً هیچ‌یک از روش‌ها سازوکاری برای \emph{کنترل} سطح تقویت خروجی ندارند و اکثراً یک نگاشت ثابت از تصویر بدون‌تزریق به یک پویش مشخص می‌آموزند. تنها استثنای مستند پینوک و همکاران~\cite{Pin23} است که یک شبکه‌ی واحد را با نشانگر پویش شرطی کردند، اما حتی این مطالعه \emph{سطح} تقویت را شرطی نکرده است. سطح تقویت واقعی به دوز، زمان‌بندی بولوس و فیزیولوژی بیمار وابسته است و یک کمیت \emph{پیوسته} است، نه یک برچسب گسسته مانند پویش.

از این مرور چهار جای خالی مکمل استخراج می‌شود:
\begin{enumerate}
\item نبود زنجیره‌ی هم‌ترازسازی‌ای که \emph{دامنه‌ی محاسبه‌ی} معیار تراز را (نه صرفاً وزن آن را) به نواحی پایدار در برابر پویش تزریق محدود کند؛
\item نبود هرگونه اندازه‌گیری و گزارش کیفیت هم‌ترازی جفت‌های آموزشی هم‌نهادی، و در نتیجه نبود آستانه‌ای که هم‌ترازی را به‌عنوان \emph{معیار پذیرش داده} به کار ببرد؛
\item نبود مدل هم‌نهاده که خروجی‌اش را بتوان \emph{هم‌زمان} بر پویش و بر سطح تقویت شرطی کرد، و نبود معیاری که این کنترل‌پذیری را بسنجد؛ شرطی‌سازی بر پویش به‌تنهایی پیشینه دارد~\cite{Pin23}، اما شرطی‌سازی بر سطح تقویتِ پیوسته، چه به‌تنهایی و چه همراه پویش، گزارش نشده است؛
\item فقدان شواهد کمّی \emph{سامانه\footnote{\lr{System}}اتیک و بازتولیدپذیر} درباره‌ی اثر هم‌نهادی بر شکاف دقت بخش‌بندی چندعضوی؛ تنها استثنای مستند \lr{CKAP-Net}~\cite{Zho24} است (دایس بخش‌بندی روی \lr{NCCT} برابر \عدد{۷۸٫۷۰}، روی \lr{CECT} هم‌نهاده \عدد{۸۱٫۲۱} و روی \lr{CECT} واقعی \عدد{۹۲٫۲۸})، اما این شواهد از یک مطالعه‌ی منفرد، روی داده‌ی خصوصی و بدون کد منتشرشده به دست آمده‌اند.
\end{enumerate}

جدول~\ref{tab:lit-competitors} مدل‌هایی را فهرست می‌کند که فصل~\ref{chap:experiments} \lr{PLC-Net} را با آن‌ها می‌سنجد. این‌ها رقبای اصلی مسئله‌ی سرتاسری‌اند و برای آن‌که مقایسه کنترل‌شده باشد، همگی روی همان تقسیم داده و با همان ابزار امتیازدهی بازآموزی می‌شوند؛ رقبای هر زیرمؤلفه (مانند روش‌های هم‌ترازسازی) جداگانه در بخش خود آزموده می‌شوند.

\begin{table}[H]
\centering
\footnotesize
\caption{رقبای اصلی هم‌نهادی سرتاسری که در فصل~\ref{chap:experiments} بازآموزی و با \lr{PLC-Net} مقایسه می‌شوند.}
\label{tab:lit-competitors}
\begin{tabular}{|p{3.4cm}|p{4.3cm}|p{5.9cm}|}
\hline
\textbf{مدل} & \textbf{راهبرد} & \textbf{نقش در مقایسه} \\
\hline
\lr{pix2pixHD}~\cite{wang2018pix2pixhd} & \lr{GAN} شرطی جفت‌شده با مولد و متمایزکننده‌ی چندمقیاسی & مبنای ترجمه‌ی جفت‌شده؛ ارزیابی تصویری و پایین‌دستی \\
\hline
\lr{CycleGAN}~\cite{zhu2017unpaired} & سازگاری چرخه‌ای، آموزش بدون‌جفت & شاهد بدون‌جفت؛ ارزیابی تصویری و پایین‌دستی \\
\hline
\lr{ResViT}~\cite{dalmaz2022resvit} & \lr{GAN} با بلوک‌های ترادیس‌ساز باقیمانده & نزدیک‌ترین مبنای منتشرشده از نظر هدف؛ ارزیابی تصویری و پایین‌دستی \\
\hline
\lr{CyTran}~\cite{Ris21} & ترادیس‌ساز سازگار چرخه‌ای & ارزیابی تصویری و پایین‌دستی \\
\hline
\lr{DDPM}~\cite{ho2020denoising} & مدل انتشار شرطی & مبنای انتشاری؛ ارزیابی تصویری و پایین‌دستی \\
\hline
\lr{SwinUNETR}~\cite{hatamizadeh2022swinunetr}، \lr{TransUNet}~\cite{chen2021transunet} & معماری‌های بخش‌بندی به‌جای ستون فقرات مولد & جداسازی اثر \emph{معماری}؛ ارزیابی تصویری و پایین‌دستی \\
\hline
\lr{Ea-GAN}~\cite{yu2019eagans} & \lr{GAN} آگاه از لبه & ارزیابی تصویری و پایین‌دستی \\
\hline
\end{tabular}
\end{table}

هیچ‌یک از این رقبا برای هم‌نهادی چندپویشی سی‌تی شکمی طراحی نشده‌اند و هیچ‌کدام سازوکاری برای شرطی‌سازی بر پویش یا سطح تقویت ندارند؛ نقطه‌ی قوت مشترکشان پیاده‌سازی باز و در دسترس است که بازآموزی و مقایسه‌ی کنترل‌شده را ممکن می‌کند، و نقطه‌ضعف مشترکشان وابستگی کامل به جفت آموزشیِ هم‌تراز (یا، برای \lr{CycleGAN}، صرف‌نظرکردن کامل از نظارت تصویردانه‌ای) و نبود هرگونه کنترل صریح روی پویش یا سطح تقویت خروجی است.

